\documentclass[../main.tex]{subfiles}
\begin{document}

\chapter{I/O 管理}
\lessoninfo{
  video={P3L5},
  textbook={OSTEP \cite{ostep} ch.~36--37，39--42；Silberschatz ら \cite{silberschatz2018} ch.~12；Kerrisk \cite{kerrisk2010} ch.~13--14},
  keywords={デバイスレジスタ，デバイスドライバ，ブロックデバイス，キャラクタデバイス，メモリマップドI/O，ポーリング，プログラム制御I/O，DMA，OSバイパス，仮想ファイルシステム，iノード，ext2，バッファキャッシュ，ジャーナリング},
}

第7章と第8章で述べた CPU とメモリの管理に加えて，オペレーティングシステムは，
計算機が入力を受け取り，出力を生み，データを永続的に格納するためのデバイスも
管理する．この章では，オペレーティングシステムがそうしたデバイスをどのように
表現し制御するか，すなわちデバイスが見せるインタフェース，デバイス固有の知識を
包み込むドライバ，そして CPU とデバイスがコマンド・データ・通知をやり取りする
方法を調べる．続いて，ファイルアクセスがブロックデバイススタックを下る経路を
たどり，ファイルシステム---とりわけ仮想ファイルシステムと ext2---がデータを
ディスク上にどう編成し，アクセスのコストをどう減らすかを見る．

\section{I/O デバイス}
\label{sec:l11-devices}

入出力（I/O）デバイスは，計算機をその環境に接続する．\source{P3L5，レッスンの
概観，I/O デバイス，I/O デバイスの基本的な特徴；OSTEP ch.~36；Silberschatz ら
ch.~12．}デバイスの機能は大きく異なり---計算機に入力を届けるだけのもの，
計算機からの出力を受け取るだけのもの，その両方を行うもの---速度も，毎秒数文字を
生むデバイスからギガビット毎秒で転送するネットワークインタフェースまで幅広い．
オペレーティングシステムの I/O 管理は，次の3つの原則に従う．
\begin{itemize}
  \item デバイスにアクセスするための標準的なインタフェースである
    \emph{プロトコル}を定める．
  \item デバイス固有の処理を担う\emph{専用のハンドラ}，すなわちデバイス
    ドライバと割り込みハンドラを提供する．
  \item 個々のデバイスの詳細をシステムの他の部分から\emph{切り離し}，
    アプリケーションがデバイスの操作ではなくファイルなどの抽象化を使えるように
    する．
\end{itemize}

\begin{definition}[デバイスレジスタ]
  \term[でばいすれじすた]{デバイスレジスタ}[device register]とは，CPU が
  デバイスとやり取りするために読み書きする，デバイス側のレジスタである．
  デバイスは通常，何をすべきかを CPU が書き込む\emph{コマンド}レジスタ，
  デバイスとの間でデータを転送する\emph{データ}レジスタ，そしてデバイスが
  何をしているかを CPU が読み取る\emph{ステータス}レジスタを備える．
\end{definition}

これらのレジスタがデバイスのインタフェースを成す（\cref{fig:l11-device}）．%
\caution{デバイスレジスタは CPU のレジスタではない．デバイスに属し，命令中の
レジスタ名ではなくアドレスで指定される．}%
その背後にあるのがデバイスの内部構成であり，デバイス自身の CPU である
マイクロコントローラ，デバイス上のメモリ，そしてマイクロフォンのアナログ・
デジタル変換器のようなデバイス固有の回路からなる．オペレーティングシステム
から見えるのはインタフェースだけであり，デバイスが内部構成をどう使って
コマンドを実行するかは隠されている．

\begin{figure}[htbp]
  \centering
  % A basic I/O device: an interface of device registers that the CPU reads
% and writes, and device-specific internals hidden behind it.
% Usage: % A basic I/O device: an interface of device registers that the CPU reads
% and writes, and device-specific internals hidden behind it.
% Usage: % A basic I/O device: an interface of device registers that the CPU reads
% and writes, and device-specific internals hidden behind it.
% Usage: \input{figures/l11-device}   (width ~116 mm)
\begin{tikzpicture}[x=1mm, y=1mm, font=\footnotesize\sffamily,
  >={Stealth[length=1.6mm]},
  reg/.style={draw, fill=white, minimum width=24mm, minimum height=6mm,
              inner sep=1pt},
  part/.style={draw, fill=white, minimum width=44mm, minimum height=6mm,
               inner sep=1pt, font=\scriptsize\sffamily},
  lab/.style={font=\scriptsize\sffamily, align=center},
  hdr/.style={font=\scriptsize\sffamily\bfseries, text=ln-accent}]
  % device outline
  \draw[fill=ln-box] (30,-13) rectangle (118,19.5);
  \node[hdr] at (50,15.5) {\iflangja{インタフェース（デバイスレジスタ）}{interface (device registers)}};
  \node[hdr] at (94,15.5) {\iflangja{内部構成}{internals}};
  \draw[densely dashed, ln-muted] (69,-10.5) -- (69,12);
  \node[reg] (st) at (50,7) {\iflangja{ステータス}{status}};
  \node[reg] (cm) at (50,0) {\iflangja{コマンド}{command}};
  \node[reg] (dt) at (50,-7) {\iflangja{データ}{data}};
  \node[part] (mc) at (94,7) {\iflangja{マイクロコントローラ（デバイスの CPU）}{microcontroller (the device's CPU)}};
  \node[part] (mm) at (94,0) {\iflangja{デバイス上のメモリ}{on-device memory}};
  \node[part] (ch) at (94,-7) {\iflangja{その他の回路（A/D 変換器など）}{other chips (e.g.\ A/D converter)}};
  \draw[<->] (cm.east) -- (mm.west);
  % CPU side
  \node[draw, fill=ln-accent!22, minimum width=12mm, minimum height=8mm]
    (cpu) at (7,0) {CPU};
  \draw[<->, thick] (cpu.east) -- (34,0);
  \node[lab, anchor=south] at (22,0.6) {\iflangja{読み書き}{read /}};
  \node[lab, anchor=north] at (22,-0.6) {\iflangja{}{write}};
  \draw (34,7) -- (34,-7);
  \foreach \y in {7,0,-7} \draw[->] (34,\y) -- (38,\y);
\end{tikzpicture}
   (width ~116 mm)
\begin{tikzpicture}[x=1mm, y=1mm, font=\footnotesize\sffamily,
  >={Stealth[length=1.6mm]},
  reg/.style={draw, fill=white, minimum width=24mm, minimum height=6mm,
              inner sep=1pt},
  part/.style={draw, fill=white, minimum width=44mm, minimum height=6mm,
               inner sep=1pt, font=\scriptsize\sffamily},
  lab/.style={font=\scriptsize\sffamily, align=center},
  hdr/.style={font=\scriptsize\sffamily\bfseries, text=ln-accent}]
  % device outline
  \draw[fill=ln-box] (30,-13) rectangle (118,19.5);
  \node[hdr] at (50,15.5) {\iflangja{インタフェース（デバイスレジスタ）}{interface (device registers)}};
  \node[hdr] at (94,15.5) {\iflangja{内部構成}{internals}};
  \draw[densely dashed, ln-muted] (69,-10.5) -- (69,12);
  \node[reg] (st) at (50,7) {\iflangja{ステータス}{status}};
  \node[reg] (cm) at (50,0) {\iflangja{コマンド}{command}};
  \node[reg] (dt) at (50,-7) {\iflangja{データ}{data}};
  \node[part] (mc) at (94,7) {\iflangja{マイクロコントローラ（デバイスの CPU）}{microcontroller (the device's CPU)}};
  \node[part] (mm) at (94,0) {\iflangja{デバイス上のメモリ}{on-device memory}};
  \node[part] (ch) at (94,-7) {\iflangja{その他の回路（A/D 変換器など）}{other chips (e.g.\ A/D converter)}};
  \draw[<->] (cm.east) -- (mm.west);
  % CPU side
  \node[draw, fill=ln-accent!22, minimum width=12mm, minimum height=8mm]
    (cpu) at (7,0) {CPU};
  \draw[<->, thick] (cpu.east) -- (34,0);
  \node[lab, anchor=south] at (22,0.6) {\iflangja{読み書き}{read /}};
  \node[lab, anchor=north] at (22,-0.6) {\iflangja{}{write}};
  \draw (34,7) -- (34,-7);
  \foreach \y in {7,0,-7} \draw[->] (34,\y) -- (38,\y);
\end{tikzpicture}
   (width ~116 mm)
\begin{tikzpicture}[x=1mm, y=1mm, font=\footnotesize\sffamily,
  >={Stealth[length=1.6mm]},
  reg/.style={draw, fill=white, minimum width=24mm, minimum height=6mm,
              inner sep=1pt},
  part/.style={draw, fill=white, minimum width=44mm, minimum height=6mm,
               inner sep=1pt, font=\scriptsize\sffamily},
  lab/.style={font=\scriptsize\sffamily, align=center},
  hdr/.style={font=\scriptsize\sffamily\bfseries, text=ln-accent}]
  % device outline
  \draw[fill=ln-box] (30,-13) rectangle (118,19.5);
  \node[hdr] at (50,15.5) {\iflangja{インタフェース（デバイスレジスタ）}{interface (device registers)}};
  \node[hdr] at (94,15.5) {\iflangja{内部構成}{internals}};
  \draw[densely dashed, ln-muted] (69,-10.5) -- (69,12);
  \node[reg] (st) at (50,7) {\iflangja{ステータス}{status}};
  \node[reg] (cm) at (50,0) {\iflangja{コマンド}{command}};
  \node[reg] (dt) at (50,-7) {\iflangja{データ}{data}};
  \node[part] (mc) at (94,7) {\iflangja{マイクロコントローラ（デバイスの CPU）}{microcontroller (the device's CPU)}};
  \node[part] (mm) at (94,0) {\iflangja{デバイス上のメモリ}{on-device memory}};
  \node[part] (ch) at (94,-7) {\iflangja{その他の回路（A/D 変換器など）}{other chips (e.g.\ A/D converter)}};
  \draw[<->] (cm.east) -- (mm.west);
  % CPU side
  \node[draw, fill=ln-accent!22, minimum width=12mm, minimum height=8mm]
    (cpu) at (7,0) {CPU};
  \draw[<->, thick] (cpu.east) -- (34,0);
  \node[lab, anchor=south] at (22,0.6) {\iflangja{読み書き}{read /}};
  \node[lab, anchor=north] at (22,-0.6) {\iflangja{}{write}};
  \draw (34,7) -- (34,-7);
  \foreach \y in {7,0,-7} \draw[->] (34,\y) -- (38,\y);
\end{tikzpicture}

  \caption{基本的な I/O デバイス．CPU はコマンド・ステータス・データの各
    レジスタを通じてのみデバイスとやり取りする．コマンドを実行する内部構成は
    デバイス固有である．}
  \label{fig:l11-device}
\end{figure}

\section{CPU とデバイスのインターコネクト}
\label{sec:l11-interconnect}

デバイスは，\source{P3L5，CPU とデバイスのインターコネクト；OSTEP ch.~36；
Silberschatz ら ch.~12．}インターコネクト上のコントローラを介して CPU と
メモリに接続される．インターコネクトは，デバイスとシステムの他の部分との間で
データをどれだけ速く動かせるか，またいくつのデバイスを接続できるかを制限する．
そうしたインターコネクトの標準的な一族が \term{PCI}[Peripheral Component Interconnect]%
\index{Peripheral Component Interconnect|see{PCI}} である．当初の形態は共有の
並列バスであり，PCI-X（PCI eXtended）がそれを高速化した．現在のプラット
フォームは PCI Express（PCIe）を用いる．PCIe は PCI のソフトウェア
インタフェースとの互換性を保ちながら，より大きな帯域幅，より速く低遅延な
アクセス，より多くのデバイスの接続を提供する．\margin{PCIe は慣習的にバスと
呼ばれるが，各デバイスはポイントツーポイントのシリアルリンクとスイッチを
通じて接続される．PCIe~4.0 の1レーンは片方向あたり約
\SI{2}{\giga\byte\per\second}，$\times 16$ のスロットは約
\SI{32}{\giga\byte\per\second} である（PCI-SIG）．}

すべてのデバイスが PCIe に直接つながるわけではない
（\cref{fig:l11-interconnect}）．PCIe に接続された SCSI コントローラは自身の
SCSI バスにディスクをつなぎ，拡張バスインタフェースはキーボードのような低速な
周辺機器を拡張（周辺）バスにつなぐ．ブリッジとなるコントローラがバスどうしの
差異を吸収し，CPU とメモリを結ぶブリッジは両者を PCIe にも接続する．したがって
CPU からデバイスへのアクセスは，複数のインターコネクトをまたぐことがある．

\begin{figure}[htbp]
  \centering
  % Typical interconnect: CPU and memory behind a bridge/memory controller,
% device controllers on PCIe, and further buses behind their controllers.
% Usage: % Typical interconnect: CPU and memory behind a bridge/memory controller,
% device controllers on PCIe, and further buses behind their controllers.
% Usage: % Typical interconnect: CPU and memory behind a bridge/memory controller,
% device controllers on PCIe, and further buses behind their controllers.
% Usage: \input{figures/l11-interconnect}   (width ~116 mm)
\begin{tikzpicture}[x=1mm, y=1mm, font=\scriptsize\sffamily,
  box/.style={draw, fill=white, minimum height=6mm, inner sep=1pt, align=center},
  ctl/.style={draw, fill=ln-box, minimum height=8.5mm, minimum width=25mm,
              inner sep=1pt, align=center},
  bus/.style={line width=1.6pt},
  blab/.style={font=\scriptsize\sffamily\itshape, text=ln-muted}]
  % top row
  \node[box, fill=ln-accent!22, minimum width=18mm, minimum height=8.5mm]
    (cpu) at (16,27) {CPU};
  \node[ctl, minimum width=34mm] (br) at (58,27)
    {\iflangja{ブリッジ／\\メモリコントローラ}{bridge /\\memory controller}};
  \node[box, minimum width=22mm, minimum height=8.5mm] (mem) at (100,27)
    {\iflangja{メモリ}{memory}};
  \draw[thick] (cpu) -- (br);
  \draw[thick] (br) -- (mem);
  % PCIe
  \draw[bus] (3,14) -- (113,14);
  \node[blab, anchor=south west] at (3,14.6) {PCIe};
  \draw[thick] (br.south) -- (58,14);
  % controllers
  \node[ctl] (gc) at (16,3) {\iflangja{グラフィックス\\コントローラ}{graphics\\controller}};
  \node[ctl] (ni) at (44,3) {\iflangja{ネットワーク\\インタフェース}{network\\interface}};
  \node[ctl] (sc) at (72,3) {\iflangja{SCSI\\コントローラ}{SCSI\\controller}};
  \node[ctl] (ex) at (100,3) {\iflangja{拡張バス\\インタフェース}{expansion bus\\interface}};
  \foreach \n in {gc,ni,sc,ex} \draw[thick] (\n.north) -- (\n.north |- 0,14);
  % below the controllers
  \node[box, minimum width=18mm] (disp) at (16,-15) {\iflangja{ディスプレイ}{display}};
  \draw (gc.south) -- (disp.north);
  \node[box, minimum width=18mm] (net) at (44,-15) {\iflangja{ネットワーク}{network}};
  \draw (ni.south) -- (net.north);
  \draw[bus] (59.5,-9) -- (84.5,-9);
  \node[blab, anchor=south west] at (59.5,-8.6) {\iflangja{SCSI バス}{SCSI bus}};
  \draw (80,-1.25) -- (80,-9);
  \node[box, minimum width=10mm] (d1) at (66,-17) {\iflangja{ディスク}{disk}};
  \node[box, minimum width=10mm] (d2) at (78,-17) {\iflangja{ディスク}{disk}};
  \draw (d1.north) -- (66,-9);  \draw (d2.north) -- (78,-9);
  \draw[bus] (87.5,-9) -- (112.5,-9);
  \node[blab, anchor=south west] at (87.5,-8.6) {\iflangja{拡張バス}{expansion bus}};
  \draw (109,-1.25) -- (109,-9);
  \node[box, minimum width=13mm] (kb) at (94,-17) {\iflangja{キーボード}{keyboard}};
  \node[box, minimum width=10mm] (ms) at (107,-17) {\iflangja{マウス}{mouse}};
  \draw (kb.north) -- (94,-9);  \draw (ms.north) -- (107,-9);
\end{tikzpicture}
   (width ~116 mm)
\begin{tikzpicture}[x=1mm, y=1mm, font=\scriptsize\sffamily,
  box/.style={draw, fill=white, minimum height=6mm, inner sep=1pt, align=center},
  ctl/.style={draw, fill=ln-box, minimum height=8.5mm, minimum width=25mm,
              inner sep=1pt, align=center},
  bus/.style={line width=1.6pt},
  blab/.style={font=\scriptsize\sffamily\itshape, text=ln-muted}]
  % top row
  \node[box, fill=ln-accent!22, minimum width=18mm, minimum height=8.5mm]
    (cpu) at (16,27) {CPU};
  \node[ctl, minimum width=34mm] (br) at (58,27)
    {\iflangja{ブリッジ／\\メモリコントローラ}{bridge /\\memory controller}};
  \node[box, minimum width=22mm, minimum height=8.5mm] (mem) at (100,27)
    {\iflangja{メモリ}{memory}};
  \draw[thick] (cpu) -- (br);
  \draw[thick] (br) -- (mem);
  % PCIe
  \draw[bus] (3,14) -- (113,14);
  \node[blab, anchor=south west] at (3,14.6) {PCIe};
  \draw[thick] (br.south) -- (58,14);
  % controllers
  \node[ctl] (gc) at (16,3) {\iflangja{グラフィックス\\コントローラ}{graphics\\controller}};
  \node[ctl] (ni) at (44,3) {\iflangja{ネットワーク\\インタフェース}{network\\interface}};
  \node[ctl] (sc) at (72,3) {\iflangja{SCSI\\コントローラ}{SCSI\\controller}};
  \node[ctl] (ex) at (100,3) {\iflangja{拡張バス\\インタフェース}{expansion bus\\interface}};
  \foreach \n in {gc,ni,sc,ex} \draw[thick] (\n.north) -- (\n.north |- 0,14);
  % below the controllers
  \node[box, minimum width=18mm] (disp) at (16,-15) {\iflangja{ディスプレイ}{display}};
  \draw (gc.south) -- (disp.north);
  \node[box, minimum width=18mm] (net) at (44,-15) {\iflangja{ネットワーク}{network}};
  \draw (ni.south) -- (net.north);
  \draw[bus] (59.5,-9) -- (84.5,-9);
  \node[blab, anchor=south west] at (59.5,-8.6) {\iflangja{SCSI バス}{SCSI bus}};
  \draw (80,-1.25) -- (80,-9);
  \node[box, minimum width=10mm] (d1) at (66,-17) {\iflangja{ディスク}{disk}};
  \node[box, minimum width=10mm] (d2) at (78,-17) {\iflangja{ディスク}{disk}};
  \draw (d1.north) -- (66,-9);  \draw (d2.north) -- (78,-9);
  \draw[bus] (87.5,-9) -- (112.5,-9);
  \node[blab, anchor=south west] at (87.5,-8.6) {\iflangja{拡張バス}{expansion bus}};
  \draw (109,-1.25) -- (109,-9);
  \node[box, minimum width=13mm] (kb) at (94,-17) {\iflangja{キーボード}{keyboard}};
  \node[box, minimum width=10mm] (ms) at (107,-17) {\iflangja{マウス}{mouse}};
  \draw (kb.north) -- (94,-9);  \draw (ms.north) -- (107,-9);
\end{tikzpicture}
   (width ~116 mm)
\begin{tikzpicture}[x=1mm, y=1mm, font=\scriptsize\sffamily,
  box/.style={draw, fill=white, minimum height=6mm, inner sep=1pt, align=center},
  ctl/.style={draw, fill=ln-box, minimum height=8.5mm, minimum width=25mm,
              inner sep=1pt, align=center},
  bus/.style={line width=1.6pt},
  blab/.style={font=\scriptsize\sffamily\itshape, text=ln-muted}]
  % top row
  \node[box, fill=ln-accent!22, minimum width=18mm, minimum height=8.5mm]
    (cpu) at (16,27) {CPU};
  \node[ctl, minimum width=34mm] (br) at (58,27)
    {\iflangja{ブリッジ／\\メモリコントローラ}{bridge /\\memory controller}};
  \node[box, minimum width=22mm, minimum height=8.5mm] (mem) at (100,27)
    {\iflangja{メモリ}{memory}};
  \draw[thick] (cpu) -- (br);
  \draw[thick] (br) -- (mem);
  % PCIe
  \draw[bus] (3,14) -- (113,14);
  \node[blab, anchor=south west] at (3,14.6) {PCIe};
  \draw[thick] (br.south) -- (58,14);
  % controllers
  \node[ctl] (gc) at (16,3) {\iflangja{グラフィックス\\コントローラ}{graphics\\controller}};
  \node[ctl] (ni) at (44,3) {\iflangja{ネットワーク\\インタフェース}{network\\interface}};
  \node[ctl] (sc) at (72,3) {\iflangja{SCSI\\コントローラ}{SCSI\\controller}};
  \node[ctl] (ex) at (100,3) {\iflangja{拡張バス\\インタフェース}{expansion bus\\interface}};
  \foreach \n in {gc,ni,sc,ex} \draw[thick] (\n.north) -- (\n.north |- 0,14);
  % below the controllers
  \node[box, minimum width=18mm] (disp) at (16,-15) {\iflangja{ディスプレイ}{display}};
  \draw (gc.south) -- (disp.north);
  \node[box, minimum width=18mm] (net) at (44,-15) {\iflangja{ネットワーク}{network}};
  \draw (ni.south) -- (net.north);
  \draw[bus] (59.5,-9) -- (84.5,-9);
  \node[blab, anchor=south west] at (59.5,-8.6) {\iflangja{SCSI バス}{SCSI bus}};
  \draw (80,-1.25) -- (80,-9);
  \node[box, minimum width=10mm] (d1) at (66,-17) {\iflangja{ディスク}{disk}};
  \node[box, minimum width=10mm] (d2) at (78,-17) {\iflangja{ディスク}{disk}};
  \draw (d1.north) -- (66,-9);  \draw (d2.north) -- (78,-9);
  \draw[bus] (87.5,-9) -- (112.5,-9);
  \node[blab, anchor=south west] at (87.5,-8.6) {\iflangja{拡張バス}{expansion bus}};
  \draw (109,-1.25) -- (109,-9);
  \node[box, minimum width=13mm] (kb) at (94,-17) {\iflangja{キーボード}{keyboard}};
  \node[box, minimum width=10mm] (ms) at (107,-17) {\iflangja{マウス}{mouse}};
  \draw (kb.north) -- (94,-9);  \draw (ms.north) -- (107,-9);
\end{tikzpicture}

  \caption{典型的なプラットフォームのインターコネクト．デバイスコントローラ
    は PCIe に接続され，ディスクとキーボードはそれぞれのコントローラの先の
    別のバスに接続される．PCIe は簡単のためバスとして描いてあるが，実際には
    各コントローラは自身のポイントツーポイントリンクで接続される．}
  \label{fig:l11-interconnect}
\end{figure}

\section{デバイスドライバ}
\label{sec:l11-drivers}

オペレーティングシステム自身が，\source{P3L5，デバイスドライバ；OSTEP
ch.~36；Silberschatz ら ch.~12．}各デバイスを操作する方法の知識を持っている
わけではない．

\begin{definition}[デバイスドライバ]
  \term[でばいすどらいば]{デバイスドライバ}[device driver]とは，ある一種類の
  デバイスに対するすべてのアクセス・管理・制御を担う，オペレーティング
  システムのデバイス固有のソフトウェア構成要素である．
\end{definition}

ドライバはデバイスの種類ごとに1つあり，通常はデバイスの製造者が提供する．
モジュール構成のオペレーティングシステムでは，カーネルモジュール（第1章）
として読み込まれる．一方でオペレーティングシステムは，ドライバが実装すべき
インタフェース，すなわち\emph{デバイスドライバフレームワーク}を標準化する．
この標準化から2つの性質が得られる．
\begin{description}
  \item[デバイス独立性] オペレーティングシステムは，ある種類のすべての
    デバイスを同じインタフェースを通じて使うので，個々のデバイスに特化する
    必要がない．
  \item[デバイスの多様性] オペレーティングシステムは，新しいものを含め，
    任意に異なるデバイスを扱える．そのデバイス用のインタフェースを実装する
    ドライバがあれば，デバイスはすぐに使える．
\end{description}

\section{デバイスの種類}
\label{sec:l11-types}

オペレーティングシステムはデバイスをいくつかのクラスに分類する．\source{P3L5，
デバイスの種類；Silberschatz ら ch.~12；Kerrisk ch.~14．}同じクラスの
デバイスは同じやり方で使われ，クラスごとに固有のドライバインタフェースを持つ．
区別されるクラスは3つである．
\begin{itemize}
  \item ディスクのような\term[ぶろっくでばいす]{ブロックデバイス}[block device]は，
    データを固定長のブロックに格納し，個々のブロックを番号で直接読み書き
    できる．
  \item キーボードのような
    \term[きゃらくたでばいす]{キャラクタデバイス}[character device]は，文字の
    逐次的なストリームを渡し，あるいは受け取る．基本操作は1文字の取得と出力で
    ある．
  \item イーサネットインタフェースのような
    \term[ねっとわーくでばいす]{ネットワークデバイス}[network device]は，
    両者の中間に位置する．データのストリームを渡すが，固定長のブロックでも
    1文字ずつでもなく，大きさの異なる塊で渡す．
\end{itemize}

\subsection{ファイルとしてのデバイス}

オペレーティングシステムは，デバイスをファイルシステムの中に表現する．

\begin{definition}[デバイスファイル]
  \term[でばいすふぁいる]{デバイスファイル}[device file]は，
  \emph{スペシャルファイル}とも呼ばれ，デバイスを表すファイルである．
  プロセスがそれに対して行う操作---open，read，write，close---は，その
  デバイスのドライバによってデバイス上で実行される．
\end{definition}

Unix 系システムでは，デバイスファイルは \texttt{/dev} ディレクトリに置かれ，
このディレクトリは devfs のような特別なファイルシステム，現在の Linux では
tmpfs に基づくものが管理する（tmpfs は第9章の POSIX 共有メモリの置き場でも
あった）．\texttt{ls -l} による詳細表示では，ブロックデバイスファイルは種別
\texttt{b}，キャラクタデバイスファイルは \texttt{c} と示される．\margin{Linux
ではネットワークデバイスが例外で，デバイスファイルを持たず，ソケットを通じて
使われる．}デバイスがファイルに見えるので，デバイスについて何も知らない
プログラムでもそれを使える．

\begin{example}
  次のコマンドは，1台目のディスクの先頭 512 バイトを通常のファイルへコピー
  する．
  \begin{lstlisting}[numbers=none]
dd if=/dev/sda of=head.img bs=512 count=1
\end{lstlisting}
  \texttt{dd} は open と read と write しか使わない．入力がブロックデバイスで
  あることも，ブロックを取り出すディスクのプロトコルも，オペレーティング
  システムとディスクのドライバが担う．
\end{example}

すべてのデバイスファイルがハードウェアに対応するわけではない．

\begin{definition}[疑似デバイス]
  \term[ぎじでばいす]{疑似デバイス}[pseudo device]とは，他と同じくデバイス
  ファイルで表現されるが，物理的なハードウェアに対応しないデバイスである．
  そのドライバは，デバイスの機能をすべてソフトウェアで実装する．
\end{definition}

例えば \texttt{/dev/zero} を読むとゼロバイトの無限のストリームが得られ，
指定した大きさのゼロで埋めたファイルを作るのに便利である．

\begin{keypoint}
  デバイスは，コマンド・データ・ステータスの各レジスタを通じて，デバイス
  ドライバによって操作される．ドライバはデバイスを，オペレーティングシステム
  がクラス全体に対して標準化したインタフェースの背後に隠す．クラスはブロック
  デバイス，キャラクタデバイス，ネットワークデバイスであり，前の2つはデバイス
  ファイルとして表現されるので，通常のファイル操作でデバイスを扱える．
\end{keypoint}

\section{CPU とデバイスの相互作用}
\label{sec:l11-interactions}

デバイスを使うには，CPU がそのレジスタをアドレス指定でき，デバイスが注意を
必要とするときに CPU へ伝えられなければならない．\source{P3L5，CPU と
デバイスの相互作用；OSTEP ch.~36；Silberschatz ら ch.~12．}レジスタを
アドレス指定する方法は2つある．

\begin{definition}[メモリマップドI/O]
  \term[めもりまっぷどIO]{メモリマップドI/O}[memory-mapped I/O]では，
  デバイスレジスタが特定の物理アドレスのメモリ位置として CPU に見える．
  物理アドレス空間の一部がデバイスとのやり取りに充てられ，CPU はそのアドレス
  に対する通常のロード命令やストア命令でレジスタにアクセスする．
  プラットフォームはそのアクセスを，メモリではなくデバイスへ導く．%
  \caution{メモリマップドレジスタは通常のメモリではない．その領域は
  キャッシュ不可で対応付けられ，ドライバは単なるポインタ参照ではなく
  \texttt{readl} や \texttt{writel} のようなアクセサを使う．}
\end{definition}

PCI デバイスでは，デバイスのコンフィギュレーション空間にあるベースアドレス
レジスタ（BAR）が，物理アドレス空間のうちそのデバイスの領域がどれだけの
大きさで，どこに位置するかを決める．BAR はブート処理の間に設定される．
メモリマップドI/O は，ファイルの内容をプロセスの仮想アドレス空間に対応付ける
メモリマップドファイル（第9章）と混同してはならない．

2つ目の方法は，x86 の \texttt{in} と \texttt{out} のような専用の CPU 命令を
使う．\term[IOぽーと]{I/Oポート}[I/O port]とは，メモリアドレスとは別に，
デバイスレジスタを指定するアドレスである．命令は対象のポートと，レジスタへ
書き込む値を指定し，あるいは読み出した値を受け取る．\margin{x86 のポート
空間は，ポート番号が16ビットであるためわずか \SI{64}{\kibi\byte} しかない．
今日のデバイスとのやり取りの大半がメモリマップドI/O で行われる理由の1つで
ある．}

\subsection{割り込みとポーリング}

逆方向では，デバイスは要求を完了したこと，あるいはデータが到着したことを CPU
に知らせなければならない．デバイスは割り込み（第5章）を発生させることができ，
あるいはオペレーティングシステムがポーリングによって知ることもできる．

\begin{definition}[ポーリング]
  デバイスが割り込みを発生させるのを待つのではなく，オペレーティングシステムが
  自ら選んだ時点でデバイスのステータスレジスタを読み，デバイスが注意を必要と
  しているかどうかを調べる技法を\term[ぽーりんぐ]{ポーリング}[polling]と
  呼ぶ．
\end{definition}

\cref{tab:l11-int-poll}に両者を比較する．割り込みはデバイスが生成すると
同時に CPU に届くが，割り込みごとに処理のオーバーヘッドがかかる．割り込み
ハンドラの実行，割り込みマスクの設定と解除に加えて，割り込まれたスレッドの
データが CPU キャッシュから追い出されるといった間接的な影響もある．
ポーリングはオペレーティングシステムにとって都合のよいときに行われるが，
まれにしか行わなければイベントは次のポーリングまで待たされ，頻繁に行えば
ほとんどのポーリングは何も見つけずに CPU 時間を浪費する．どちらが望ましいか
は，デバイスの種類と，イベントがどれだけ頻繁でどれだけ速く処理されねば
ならないかといった目的に依存する．\sidenote{Linux のネットワークドライバは
NAPI によって両者を切り替える．最初のパケットが届くとドライバはその
インタフェースの受信割り込みを禁止し，カーネルは滞留した分がなくなるまで
デバイスをポーリングし，その後で割り込みを再び許可する．負荷の軽い
インタフェースは割り込み駆動に，忙しいインタフェースはポーリングになる．}

\begin{table}[htbp]
  \centering
  \caption{割り込みとポーリング．}
  \label{tab:l11-int-poll}
  \small
  \begin{tabular}{@{}>{\raggedright\arraybackslash}p{0.2\linewidth}>{\raggedright\arraybackslash}p{0.36\linewidth}>{\raggedright\arraybackslash}p{0.36\linewidth}@{}}
    \toprule
    & 割り込み & ポーリング \\
    \midrule
    開始する主体 & デバイス & オペレーティングシステム \\
    イベントの認識 & 生成されると同時 & 次のポーリング時 \\
    コスト & 割り込みごとの処理オーバーヘッド，キャッシュ汚染 & ポーリング
      ごとに1回のステータス読み出し．何も見つからない回も含む \\
    欠点 & イベントが頻繁なときのオーバーヘッド & ポーリングがまれなら遅延，
      頻繁なら CPU のオーバーヘッド \\
    \bottomrule
  \end{tabular}
\end{table}

\begin{example}
  割り込みを受けるコストが，イベント自体の処理を除いて CPU 時間で
  \SI{4}{\micro\second} であり，前回のポーリング以降に溜まったイベントを
  すべて回収する1回のポーリングのコストが \SI{0.5}{\micro\second} であると
  する．毎秒 50 件のイベントを報告するデバイスは，割り込みでは毎秒
  \SI{200}{\micro\second} の CPU 時間を要する．\SI{1}{\milli\second} ごとの
  ポーリングなら毎秒 \SI{500}{\micro\second} を要し，各イベントは最大
  \SI{1}{\milli\second} 遅れる．一方，毎秒 \num{100000} パケットを受信する
  ネットワークインタフェースは，パケットごとに1回の割り込みでは毎秒
  \SI{400}{\milli\second}（CPU の \SI{40}{\percent}）を要するのに対し，
  \SI{1}{\milli\second} ごとのポーリングなら依然として毎秒
  \SI{500}{\micro\second} で済む．
\end{example}

\section{デバイスアクセス：プログラム制御I/O と DMA}
\label{sec:l11-pio-dma}

CPU とデバイスがデータをやり取りする方法は2つあり，データが CPU を通るか
どうかが異なる．\source{P3L5，デバイスアクセス：PIO，デバイスアクセス：DMA；
OSTEP ch.~36．}どちらの方法も，コマンドの発行には \cref{sec:l11-devices}の
デバイスレジスタを使う．

\begin{definition}[プログラム制御I/O]
  \term[ぷろぐらむせいぎょIO]{プログラム制御I/O}[programmed I/O]%
  \index{PIO|see{プログラム制御I/O}}（PIO）では，CPU がやり取りのすべての
  段階をデバイスレジスタを通じて行う．操作を要求するためにコマンドレジスタへ
  書き込み，データレジスタへの書き込みまたは読み出しによって，レジスタの
  大きさずつデータを運ぶ．
\end{definition}

PIO はレジスタ以外のハードウェアによる支援を必要としない．PIO で
ネットワークパケットを送るには，CPU は送信を要求するコマンドをコマンド
レジスタに書き，パケットの先頭部分をデータレジスタにコピーし，パケット全体が
転送されるまでコピーを繰り返す（\cref{fig:l11-pio-dma}a）．CPU のアクセス
回数はデータの大きさとともに増える．

\begin{definition}[ダイレクトメモリアクセス]
  \term{DMA}[direct memory access]\index{direct memory access|see{DMA}}では，
  CPU は依然としてコマンドレジスタを通じてコマンドを発行するが，データは CPU
  を通らない．CPU は主メモリ上のバッファにデータを置き，そのバッファの
  アドレスと大きさをデバイスに渡す．すると DMA コントローラが，メモリと
  デバイスの間でデータを自ら転送する．
\end{definition}

DMA でパケットを送るには，CPU はコマンドを書き，パケットをメモリのバッファに
置き，バッファのアドレスと大きさを DMA コントローラに設定する
（\cref{fig:l11-pio-dma}b）．CPU のアクセス回数はもはやデータの大きさに
依存しない．

\begin{figure}[htbp]
  \centering
  % Sending a packet with programmed I/O (data copied by the CPU into the data
% register) and with DMA (device copies the data from a memory buffer).
% Usage: % Sending a packet with programmed I/O (data copied by the CPU into the data
% register) and with DMA (device copies the data from a memory buffer).
% Usage: % Sending a packet with programmed I/O (data copied by the CPU into the data
% register) and with DMA (device copies the data from a memory buffer).
% Usage: \input{figures/l11-pio-dma}   (width ~118 mm)
\begin{tikzpicture}[x=1mm, y=1mm, font=\footnotesize\sffamily,
  >={Stealth[length=1.6mm]},
  box/.style={draw, fill=white, minimum height=7mm, inner sep=1pt, align=center},
  reg/.style={draw, fill=white, minimum width=20mm, minimum height=5.5mm,
              inner sep=1pt, font=\scriptsize\sffamily},
  num/.style={circle, draw=black, fill=white, text=black, line width=0.4pt,
              inner sep=0.4pt, font=\scriptsize\sffamily\bfseries},
  lab/.style={font=\scriptsize\sffamily, align=left},
  ttl/.style={font=\footnotesize\sffamily\bfseries, anchor=west}]
  % ---------------- (a) PIO
  \node[ttl] at (0,26) {\iflangja{(a) プログラム制御 I/O}{(a) programmed I/O}};
  \node[box, fill=ln-accent!22, minimum width=14mm] (cpu) at (9,13) {CPU};
  \node[box, minimum width=16mm, minimum height=9mm] (mem) at (9,-4)
    {\iflangja{メモリ\\（パケット）}{memory\\(packet)}};
  \draw[fill=ln-box] (30,-9) rectangle (56,21);
  \node[font=\scriptsize\sffamily\bfseries, text=ln-accent, anchor=north]
    at (43,20.5) {NIC};
  \node[reg] (cm) at (43,13) {\iflangja{コマンド}{command}};
  \node[reg] (dt) at (43,3) {\iflangja{データ}{data}};
  \draw[->] (cpu.east) -- node[num, above=0.3mm] {1} (cm.west);
  \draw[->, line width=1.4pt, ln-accent] (mem.north) -- (cpu.south);
  \draw[->, line width=1.4pt, ln-accent] (cpu.south east) -- node[num, below left=-0.4mm] {2} (dt.west);
  \node[lab, anchor=north west, text width=56mm] at (0,-12)
    {\iflangja{\textbf{1} コマンドレジスタへの書き込み 1 回\\
      \textbf{2} データレジスタへの書き込み（パケット全体を\\
      \phantom{\textbf{2} }レジスタの大きさずつ繰り返す）}%
     {\textbf{1} one write to the command register\\
      \textbf{2} writes to the data register, repeated\\
      \phantom{\textbf{2} }until the whole packet is transferred}};
  % ---------------- (b) DMA
  \begin{scope}[xshift=62mm]
  \node[ttl] at (0,26) {\iflangja{(b) DMA}{(b) direct memory access}};
  \node[box, fill=ln-accent!22, minimum width=14mm] (cpu2) at (9,13) {CPU};
  \node[box, minimum width=16mm, minimum height=9mm] (mem2) at (9,-4)
    {\iflangja{メモリ\\（バッファ）}{memory\\(buffer)}};
  \draw[fill=ln-box] (30,-9) rectangle (56,21);
  \node[font=\scriptsize\sffamily\bfseries, text=ln-accent, anchor=north]
    at (43,20.5) {NIC};
  \node[reg] (cm2) at (43,13) {\iflangja{コマンド}{command}};
  \node[reg, fill=ln-accent!10] (dma) at (43,-4) {\iflangja{DMA コントローラ}{DMA controller}};
  \draw[->] (cpu2.east) -- node[num, above=0.3mm] {1} (cm2.west);
  \draw[->] (cpu2.south east) -- node[num, above right=-0.6mm] {2} (dma.north west);
  \draw[->, line width=1.4pt, ln-accent] (mem2.east |- dma.west) -- node[num, below=0.8mm] {3} (dma.west);
  \node[lab, anchor=north west, text width=56mm] at (0,-12)
    {\iflangja{\textbf{1} コマンドレジスタへの書き込み 1 回\\
      \textbf{2} DMA の設定（バッファのアドレスとサイズ）\\
      \textbf{3} デバイスがバッファからデータを読む}%
     {\textbf{1} one write to the command register\\
      \textbf{2} DMA configuration: buffer address and size\\
      \textbf{3} the device reads the data from the buffer}};
  \end{scope}
\end{tikzpicture}
   (width ~118 mm)
\begin{tikzpicture}[x=1mm, y=1mm, font=\footnotesize\sffamily,
  >={Stealth[length=1.6mm]},
  box/.style={draw, fill=white, minimum height=7mm, inner sep=1pt, align=center},
  reg/.style={draw, fill=white, minimum width=20mm, minimum height=5.5mm,
              inner sep=1pt, font=\scriptsize\sffamily},
  num/.style={circle, draw=black, fill=white, text=black, line width=0.4pt,
              inner sep=0.4pt, font=\scriptsize\sffamily\bfseries},
  lab/.style={font=\scriptsize\sffamily, align=left},
  ttl/.style={font=\footnotesize\sffamily\bfseries, anchor=west}]
  % ---------------- (a) PIO
  \node[ttl] at (0,26) {\iflangja{(a) プログラム制御 I/O}{(a) programmed I/O}};
  \node[box, fill=ln-accent!22, minimum width=14mm] (cpu) at (9,13) {CPU};
  \node[box, minimum width=16mm, minimum height=9mm] (mem) at (9,-4)
    {\iflangja{メモリ\\（パケット）}{memory\\(packet)}};
  \draw[fill=ln-box] (30,-9) rectangle (56,21);
  \node[font=\scriptsize\sffamily\bfseries, text=ln-accent, anchor=north]
    at (43,20.5) {NIC};
  \node[reg] (cm) at (43,13) {\iflangja{コマンド}{command}};
  \node[reg] (dt) at (43,3) {\iflangja{データ}{data}};
  \draw[->] (cpu.east) -- node[num, above=0.3mm] {1} (cm.west);
  \draw[->, line width=1.4pt, ln-accent] (mem.north) -- (cpu.south);
  \draw[->, line width=1.4pt, ln-accent] (cpu.south east) -- node[num, below left=-0.4mm] {2} (dt.west);
  \node[lab, anchor=north west, text width=56mm] at (0,-12)
    {\iflangja{\textbf{1} コマンドレジスタへの書き込み 1 回\\
      \textbf{2} データレジスタへの書き込み（パケット全体を\\
      \phantom{\textbf{2} }レジスタの大きさずつ繰り返す）}%
     {\textbf{1} one write to the command register\\
      \textbf{2} writes to the data register, repeated\\
      \phantom{\textbf{2} }until the whole packet is transferred}};
  % ---------------- (b) DMA
  \begin{scope}[xshift=62mm]
  \node[ttl] at (0,26) {\iflangja{(b) DMA}{(b) direct memory access}};
  \node[box, fill=ln-accent!22, minimum width=14mm] (cpu2) at (9,13) {CPU};
  \node[box, minimum width=16mm, minimum height=9mm] (mem2) at (9,-4)
    {\iflangja{メモリ\\（バッファ）}{memory\\(buffer)}};
  \draw[fill=ln-box] (30,-9) rectangle (56,21);
  \node[font=\scriptsize\sffamily\bfseries, text=ln-accent, anchor=north]
    at (43,20.5) {NIC};
  \node[reg] (cm2) at (43,13) {\iflangja{コマンド}{command}};
  \node[reg, fill=ln-accent!10] (dma) at (43,-4) {\iflangja{DMA コントローラ}{DMA controller}};
  \draw[->] (cpu2.east) -- node[num, above=0.3mm] {1} (cm2.west);
  \draw[->] (cpu2.south east) -- node[num, above right=-0.6mm] {2} (dma.north west);
  \draw[->, line width=1.4pt, ln-accent] (mem2.east |- dma.west) -- node[num, below=0.8mm] {3} (dma.west);
  \node[lab, anchor=north west, text width=56mm] at (0,-12)
    {\iflangja{\textbf{1} コマンドレジスタへの書き込み 1 回\\
      \textbf{2} DMA の設定（バッファのアドレスとサイズ）\\
      \textbf{3} デバイスがバッファからデータを読む}%
     {\textbf{1} one write to the command register\\
      \textbf{2} DMA configuration: buffer address and size\\
      \textbf{3} the device reads the data from the buffer}};
  \end{scope}
\end{tikzpicture}
   (width ~118 mm)
\begin{tikzpicture}[x=1mm, y=1mm, font=\footnotesize\sffamily,
  >={Stealth[length=1.6mm]},
  box/.style={draw, fill=white, minimum height=7mm, inner sep=1pt, align=center},
  reg/.style={draw, fill=white, minimum width=20mm, minimum height=5.5mm,
              inner sep=1pt, font=\scriptsize\sffamily},
  num/.style={circle, draw=black, fill=white, text=black, line width=0.4pt,
              inner sep=0.4pt, font=\scriptsize\sffamily\bfseries},
  lab/.style={font=\scriptsize\sffamily, align=left},
  ttl/.style={font=\footnotesize\sffamily\bfseries, anchor=west}]
  % ---------------- (a) PIO
  \node[ttl] at (0,26) {\iflangja{(a) プログラム制御 I/O}{(a) programmed I/O}};
  \node[box, fill=ln-accent!22, minimum width=14mm] (cpu) at (9,13) {CPU};
  \node[box, minimum width=16mm, minimum height=9mm] (mem) at (9,-4)
    {\iflangja{メモリ\\（パケット）}{memory\\(packet)}};
  \draw[fill=ln-box] (30,-9) rectangle (56,21);
  \node[font=\scriptsize\sffamily\bfseries, text=ln-accent, anchor=north]
    at (43,20.5) {NIC};
  \node[reg] (cm) at (43,13) {\iflangja{コマンド}{command}};
  \node[reg] (dt) at (43,3) {\iflangja{データ}{data}};
  \draw[->] (cpu.east) -- node[num, above=0.3mm] {1} (cm.west);
  \draw[->, line width=1.4pt, ln-accent] (mem.north) -- (cpu.south);
  \draw[->, line width=1.4pt, ln-accent] (cpu.south east) -- node[num, below left=-0.4mm] {2} (dt.west);
  \node[lab, anchor=north west, text width=56mm] at (0,-12)
    {\iflangja{\textbf{1} コマンドレジスタへの書き込み 1 回\\
      \textbf{2} データレジスタへの書き込み（パケット全体を\\
      \phantom{\textbf{2} }レジスタの大きさずつ繰り返す）}%
     {\textbf{1} one write to the command register\\
      \textbf{2} writes to the data register, repeated\\
      \phantom{\textbf{2} }until the whole packet is transferred}};
  % ---------------- (b) DMA
  \begin{scope}[xshift=62mm]
  \node[ttl] at (0,26) {\iflangja{(b) DMA}{(b) direct memory access}};
  \node[box, fill=ln-accent!22, minimum width=14mm] (cpu2) at (9,13) {CPU};
  \node[box, minimum width=16mm, minimum height=9mm] (mem2) at (9,-4)
    {\iflangja{メモリ\\（バッファ）}{memory\\(buffer)}};
  \draw[fill=ln-box] (30,-9) rectangle (56,21);
  \node[font=\scriptsize\sffamily\bfseries, text=ln-accent, anchor=north]
    at (43,20.5) {NIC};
  \node[reg] (cm2) at (43,13) {\iflangja{コマンド}{command}};
  \node[reg, fill=ln-accent!10] (dma) at (43,-4) {\iflangja{DMA コントローラ}{DMA controller}};
  \draw[->] (cpu2.east) -- node[num, above=0.3mm] {1} (cm2.west);
  \draw[->] (cpu2.south east) -- node[num, above right=-0.6mm] {2} (dma.north west);
  \draw[->, line width=1.4pt, ln-accent] (mem2.east |- dma.west) -- node[num, below=0.8mm] {3} (dma.west);
  \node[lab, anchor=north west, text width=56mm] at (0,-12)
    {\iflangja{\textbf{1} コマンドレジスタへの書き込み 1 回\\
      \textbf{2} DMA の設定（バッファのアドレスとサイズ）\\
      \textbf{3} デバイスがバッファからデータを読む}%
     {\textbf{1} one write to the command register\\
      \textbf{2} DMA configuration: buffer address and size\\
      \textbf{3} the device reads the data from the buffer}};
  \end{scope}
\end{tikzpicture}

  \caption{パケットの送信．(a) プログラム制御I/O では，CPU がパケットを少し
    ずつデータレジスタへコピーする．(b) DMA では，CPU はコマンドの書き込みと
    DMA の設定を行うだけで，デバイスがメモリからパケットを読む．}
  \label{fig:l11-pio-dma}
\end{figure}

\begin{example}
  8 バイトのデータレジスタを持つネットワークインタフェースが 9000 バイトの
  フレームを送る．PIO では，CPU はコマンドのために1回，データのために
  $9000/8 = 1125$ 回，合計 1126 回のアクセスを行う．DMA では2回，すなわち
  コマンドに1回と DMA の設定に1回である．
\end{example}

DMA は必要な段階が少ないが，DMA の設定は1回のレジスタ書き込みよりも複雑な
操作である．そのため小さな転送では，PIO による数回のレジスタ書き込みの方が
1回の DMA 設定より安く済むことがある．PIO が必要とするデータアクセスの
コストが DMA の設定より小さい限り，PIO の方が安い．\tip{DMA が引き合うのは，
転送がレジスタ幅の何倍にも及ぶときである．数バイトの転送では PIO の方が
安い．}さらに，デバイスはバッファを物理アドレスでアクセスするため，別の要件
が生じる．バッファは転送が完了するまで物理メモリに留まらなければならない．
そのためバッファのページはピン留め（第8章）\index{ぴんどめぺーじ@ピン留めページ}
される，すなわちスワッピングの対象から外される．さもなければ，デバイスは
バッファを置き換えたページを読んだり上書きしたりしかねない．\margin{物理
アドレスをそのまま渡すことは，メモリ全体をデバイスに委ねることでもある．
そのため現在のプラットフォームは IOMMU（Intel VT-d，AMD-Vi）を介在させ，
ドライバはアドレスを渡す代わりにバッファを DMA 用に対応付ける．}

\begin{keypoint}
  CPU はメモリマップドI/O か I/Oポートを通じてデバイスレジスタにアクセスし，
  デバイスは割り込みによって CPU に知らせる．あるいはオペレーティングシステム
  がデバイスをポーリングする．データは CPU を経由してレジスタ単位で運ばれる
  （PIO）か，メモリバッファとデバイスの間で直接運ばれる（DMA）．後者は CPU
  のアクセス回数を減らすが，設定がより複雑になり，メモリのピン留めを要する．
\end{keypoint}

\section{典型的なデバイスアクセス}
\label{sec:l11-access}

アプリケーションが自らデバイスとやり取りすることは通常なく，その要求は
オペレーティングシステムのいくつもの層を通る．\source{P3L5，典型的な
デバイスアクセス，OS バイパス；OSTEP ch.~36．}デバイスアクセスの典型的な経路
は次のとおりである（\cref{fig:l11-access}a）．
\begin{enumerate}
  \item プロセスが，ソケットへのデータ送信やファイルからの読み出しのような
    操作を要求するシステムコールを行う．
  \item カーネルが，そのデバイスに対応するカーネル内スタックを実行する．
    TCP/IP スタックがパケットを組み立てる，あるいはファイルシステムが要求
    されたデータを保持するブロックを決定する．
  \item カーネルが適切なデバイスドライバ，ここではネットワークドライバか
    ブロックデバイスドライバを呼び出す．
  \item ドライバが，PIO によるレジスタ書き込みか DMA の設定によって，要求を
    デバイス上に設定する．
  \item デバイスが要求を実行する．パケットを送信する，あるいはディスクから
    ブロックを読む．
\end{enumerate}
結果と完了通知は同じ経路を逆にたどる．デバイスが割り込みを発生させるか
ポーリングされ，ドライバがデータをスタックの上へ渡し，システムコールが
プロセスへ戻る．

\begin{figure}[htbp]
  \centering
  % (a) Typical device access through the kernel; (b) operating system bypass
% with a user-level driver.
% Usage: % (a) Typical device access through the kernel; (b) operating system bypass
% with a user-level driver.
% Usage: % (a) Typical device access through the kernel; (b) operating system bypass
% with a user-level driver.
% Usage: \input{figures/l11-access}   (width ~118 mm)
\begin{tikzpicture}[x=1mm, y=1mm, font=\scriptsize\sffamily,
  >={Stealth[length=1.5mm]},
  box/.style={draw, fill=white, minimum height=7mm, inner sep=1pt, align=center},
  kbox/.style={draw, fill=ln-box, minimum height=7mm, inner sep=1pt, align=center},
  dev/.style={draw, fill=ln-accent!22, minimum height=7mm, inner sep=1pt, align=center},
  lab/.style={font=\scriptsize\sffamily, align=center},
  side/.style={font=\scriptsize\sffamily\itshape, text=ln-muted},
  ttl/.style={font=\footnotesize\sffamily\bfseries, anchor=west}]
  % boundaries
  \draw[densely dashed, ln-muted] (0,23) -- (118,23);
  \draw[densely dashed, ln-muted] (0,-11) -- (118,-11);
  \node[side, anchor=south west] at (0,23.3) {\iflangja{ユーザ}{user}};
  \node[side, anchor=north west] at (0,22.7) {\iflangja{カーネル}{kernel}};
  \node[side, anchor=north west] at (0,-11.3) {\iflangja{デバイス}{device}};
  % ---------------- (a) typical access
  \node[ttl] at (12,40) {\iflangja{(a) カーネルを経由するアクセス}{(a) access through the kernel}};
  \node[box, minimum width=34mm] (p) at (34,31) {\iflangja{ユーザプロセス}{user process}};
  \node[kbox, minimum width=34mm] (st) at (34,11)
    {\iflangja{カーネル内スタック\\（TCP/IP，ファイルシステム）}{in-kernel stack\\(TCP/IP, file system)}};
  \node[kbox, minimum width=34mm] (dr) at (34,-2) {\iflangja{デバイスドライバ}{device driver}};
  \node[dev, minimum width=34mm] (d) at (34,-21) {\iflangja{デバイス}{device}};
  \draw[->, thick] ([xshift=-6mm]p.south) -- node[lab, left, pos=0.72] {\iflangja{システムコール}{system call}} ([xshift=-6mm]st.north);
  \draw[->, thick] ([xshift=-6mm]st.south) -- ([xshift=-6mm]dr.north);
  \draw[->, thick] ([xshift=-6mm]dr.south) -- node[lab, left, pos=0.22] {PIO/DMA} ([xshift=-6mm]d.north);
  \draw[->] ([xshift=6mm]d.north) -- node[lab, right, pos=0.78] {\iflangja{割り込み}{interrupt}} ([xshift=6mm]dr.south);
  \draw[->] ([xshift=6mm]dr.north) -- ([xshift=6mm]st.south);
  \draw[->] ([xshift=6mm]st.north) -- node[lab, right, pos=0.28] {\iflangja{結果}{result}} ([xshift=6mm]p.south);
  % ---------------- (b) OS bypass
  \node[ttl] at (66,40) {\iflangja{(b) OS バイパス}{(b) operating system bypass}};
  \draw[fill=white] (66,25.5) rectangle (100,37);
  \node[lab] at (83,34.5) {\iflangja{ユーザプロセス}{user process}};
  \node[kbox, minimum width=30mm, minimum height=5mm] (ud) at (83,29.5)
    {\iflangja{ユーザレベルドライバ}{user-level driver}};
  \node[kbox, minimum width=26mm] (kd) at (104,5)
    {\iflangja{カーネルの\\ドライバ}{kernel driver}};
  \node[dev, minimum width=46mm] (d2) at (90,-21) {\iflangja{デバイス}{device}};
  \draw[<->, line width=1.4pt, ln-accent] (77,27) -- node[lab, fill=white, text=black, inner sep=0.8pt]
    {\iflangja{レジスタと\\メモリに\\直接アクセス}{direct access\\to registers\\and memory}} (77,-17.5);
  \draw[->, densely dashed] (kd.south) -- node[lab, right, pos=0.3] {\iflangja{設定}{configure}} (kd.south |- d2.north);
  \draw[->, densely dashed] (kd.north) |- node[lab, pos=0.15, right] {\iflangja{許可}{grant}} (100,31);
\end{tikzpicture}
   (width ~118 mm)
\begin{tikzpicture}[x=1mm, y=1mm, font=\scriptsize\sffamily,
  >={Stealth[length=1.5mm]},
  box/.style={draw, fill=white, minimum height=7mm, inner sep=1pt, align=center},
  kbox/.style={draw, fill=ln-box, minimum height=7mm, inner sep=1pt, align=center},
  dev/.style={draw, fill=ln-accent!22, minimum height=7mm, inner sep=1pt, align=center},
  lab/.style={font=\scriptsize\sffamily, align=center},
  side/.style={font=\scriptsize\sffamily\itshape, text=ln-muted},
  ttl/.style={font=\footnotesize\sffamily\bfseries, anchor=west}]
  % boundaries
  \draw[densely dashed, ln-muted] (0,23) -- (118,23);
  \draw[densely dashed, ln-muted] (0,-11) -- (118,-11);
  \node[side, anchor=south west] at (0,23.3) {\iflangja{ユーザ}{user}};
  \node[side, anchor=north west] at (0,22.7) {\iflangja{カーネル}{kernel}};
  \node[side, anchor=north west] at (0,-11.3) {\iflangja{デバイス}{device}};
  % ---------------- (a) typical access
  \node[ttl] at (12,40) {\iflangja{(a) カーネルを経由するアクセス}{(a) access through the kernel}};
  \node[box, minimum width=34mm] (p) at (34,31) {\iflangja{ユーザプロセス}{user process}};
  \node[kbox, minimum width=34mm] (st) at (34,11)
    {\iflangja{カーネル内スタック\\（TCP/IP，ファイルシステム）}{in-kernel stack\\(TCP/IP, file system)}};
  \node[kbox, minimum width=34mm] (dr) at (34,-2) {\iflangja{デバイスドライバ}{device driver}};
  \node[dev, minimum width=34mm] (d) at (34,-21) {\iflangja{デバイス}{device}};
  \draw[->, thick] ([xshift=-6mm]p.south) -- node[lab, left, pos=0.72] {\iflangja{システムコール}{system call}} ([xshift=-6mm]st.north);
  \draw[->, thick] ([xshift=-6mm]st.south) -- ([xshift=-6mm]dr.north);
  \draw[->, thick] ([xshift=-6mm]dr.south) -- node[lab, left, pos=0.22] {PIO/DMA} ([xshift=-6mm]d.north);
  \draw[->] ([xshift=6mm]d.north) -- node[lab, right, pos=0.78] {\iflangja{割り込み}{interrupt}} ([xshift=6mm]dr.south);
  \draw[->] ([xshift=6mm]dr.north) -- ([xshift=6mm]st.south);
  \draw[->] ([xshift=6mm]st.north) -- node[lab, right, pos=0.28] {\iflangja{結果}{result}} ([xshift=6mm]p.south);
  % ---------------- (b) OS bypass
  \node[ttl] at (66,40) {\iflangja{(b) OS バイパス}{(b) operating system bypass}};
  \draw[fill=white] (66,25.5) rectangle (100,37);
  \node[lab] at (83,34.5) {\iflangja{ユーザプロセス}{user process}};
  \node[kbox, minimum width=30mm, minimum height=5mm] (ud) at (83,29.5)
    {\iflangja{ユーザレベルドライバ}{user-level driver}};
  \node[kbox, minimum width=26mm] (kd) at (104,5)
    {\iflangja{カーネルの\\ドライバ}{kernel driver}};
  \node[dev, minimum width=46mm] (d2) at (90,-21) {\iflangja{デバイス}{device}};
  \draw[<->, line width=1.4pt, ln-accent] (77,27) -- node[lab, fill=white, text=black, inner sep=0.8pt]
    {\iflangja{レジスタと\\メモリに\\直接アクセス}{direct access\\to registers\\and memory}} (77,-17.5);
  \draw[->, densely dashed] (kd.south) -- node[lab, right, pos=0.3] {\iflangja{設定}{configure}} (kd.south |- d2.north);
  \draw[->, densely dashed] (kd.north) |- node[lab, pos=0.15, right] {\iflangja{許可}{grant}} (100,31);
\end{tikzpicture}
   (width ~118 mm)
\begin{tikzpicture}[x=1mm, y=1mm, font=\scriptsize\sffamily,
  >={Stealth[length=1.5mm]},
  box/.style={draw, fill=white, minimum height=7mm, inner sep=1pt, align=center},
  kbox/.style={draw, fill=ln-box, minimum height=7mm, inner sep=1pt, align=center},
  dev/.style={draw, fill=ln-accent!22, minimum height=7mm, inner sep=1pt, align=center},
  lab/.style={font=\scriptsize\sffamily, align=center},
  side/.style={font=\scriptsize\sffamily\itshape, text=ln-muted},
  ttl/.style={font=\footnotesize\sffamily\bfseries, anchor=west}]
  % boundaries
  \draw[densely dashed, ln-muted] (0,23) -- (118,23);
  \draw[densely dashed, ln-muted] (0,-11) -- (118,-11);
  \node[side, anchor=south west] at (0,23.3) {\iflangja{ユーザ}{user}};
  \node[side, anchor=north west] at (0,22.7) {\iflangja{カーネル}{kernel}};
  \node[side, anchor=north west] at (0,-11.3) {\iflangja{デバイス}{device}};
  % ---------------- (a) typical access
  \node[ttl] at (12,40) {\iflangja{(a) カーネルを経由するアクセス}{(a) access through the kernel}};
  \node[box, minimum width=34mm] (p) at (34,31) {\iflangja{ユーザプロセス}{user process}};
  \node[kbox, minimum width=34mm] (st) at (34,11)
    {\iflangja{カーネル内スタック\\（TCP/IP，ファイルシステム）}{in-kernel stack\\(TCP/IP, file system)}};
  \node[kbox, minimum width=34mm] (dr) at (34,-2) {\iflangja{デバイスドライバ}{device driver}};
  \node[dev, minimum width=34mm] (d) at (34,-21) {\iflangja{デバイス}{device}};
  \draw[->, thick] ([xshift=-6mm]p.south) -- node[lab, left, pos=0.72] {\iflangja{システムコール}{system call}} ([xshift=-6mm]st.north);
  \draw[->, thick] ([xshift=-6mm]st.south) -- ([xshift=-6mm]dr.north);
  \draw[->, thick] ([xshift=-6mm]dr.south) -- node[lab, left, pos=0.22] {PIO/DMA} ([xshift=-6mm]d.north);
  \draw[->] ([xshift=6mm]d.north) -- node[lab, right, pos=0.78] {\iflangja{割り込み}{interrupt}} ([xshift=6mm]dr.south);
  \draw[->] ([xshift=6mm]dr.north) -- ([xshift=6mm]st.south);
  \draw[->] ([xshift=6mm]st.north) -- node[lab, right, pos=0.28] {\iflangja{結果}{result}} ([xshift=6mm]p.south);
  % ---------------- (b) OS bypass
  \node[ttl] at (66,40) {\iflangja{(b) OS バイパス}{(b) operating system bypass}};
  \draw[fill=white] (66,25.5) rectangle (100,37);
  \node[lab] at (83,34.5) {\iflangja{ユーザプロセス}{user process}};
  \node[kbox, minimum width=30mm, minimum height=5mm] (ud) at (83,29.5)
    {\iflangja{ユーザレベルドライバ}{user-level driver}};
  \node[kbox, minimum width=26mm] (kd) at (104,5)
    {\iflangja{カーネルの\\ドライバ}{kernel driver}};
  \node[dev, minimum width=46mm] (d2) at (90,-21) {\iflangja{デバイス}{device}};
  \draw[<->, line width=1.4pt, ln-accent] (77,27) -- node[lab, fill=white, text=black, inner sep=0.8pt]
    {\iflangja{レジスタと\\メモリに\\直接アクセス}{direct access\\to registers\\and memory}} (77,-17.5);
  \draw[->, densely dashed] (kd.south) -- node[lab, right, pos=0.3] {\iflangja{設定}{configure}} (kd.south |- d2.north);
  \draw[->, densely dashed] (kd.north) |- node[lab, pos=0.15, right] {\iflangja{許可}{grant}} (100,31);
\end{tikzpicture}

  \caption{(a) 典型的なデバイスアクセスは，システムコール，カーネル内スタック，
    デバイスドライバを通る．(b) OS バイパスでは，プロセスはユーザレベル
    ドライバを通じてデバイスにアクセスし，カーネルはアクセスの設定と制御を
    行うだけである．}
  \label{fig:l11-access}
\end{figure}

\subsection{OS バイパス}

カーネルを経由することは必須ではない．

\begin{definition}[OS バイパス]
  デバイスのレジスタとメモリをユーザプロセスから直接アクセスできるように
  するデバイスアクセスの方式を
  \term[OSばいぱす]{OSバイパス}[operating system bypass]と呼ぶ．
  オペレーティングシステムはアクセスを設定した後は関与せず，プロセスは
  カーネルのコードを実行せずにデバイス操作を行う（\cref{fig:l11-access}b）．
\end{definition}

本来カーネルのドライバで実行されるコードは，\emph{ユーザレベルドライバ}，
すなわちアプリケーションにリンクされるライブラリが提供する．オペレーティング
システムは粗い粒度の制御を保持する．例えばデバイスの有効化と無効化や，別の
プロセスにアクセス権を与えることができる．\sidenote{カーネルバイパスは高性能
ネットワークでは日常的である．InfiniBand や RoCE のアダプタは verbs API を
通じてキューペアをプロセスに見せ，DPDK はイーサネットアダプタを完全に
ユーザ空間から駆動する．カーネルは設定を行うだけである．第15章はこうした
アダプタをリモートダイレクトメモリアクセスに使う．}バイパスはデバイスの
機能に依存する．
\begin{itemize}
  \item デバイスは十分な数のレジスタを備えていなければならない．一部を
    ユーザプロセスに対応付けつつ，制御に必要なレジスタはカーネルだけが
    アクセスできるように残すためである．
  \item 複数のプロセスがデバイスを使う場合，デバイスは\emph{逆多重化}，
    すなわちデータがどのプロセスのものかを判別できなければならない．典型的な
    デバイスアクセスではカーネルがこれを行う．例えばネットワークスタックは，
    到着したパケットのヘッダを調べて受信プロセスを見つける．バイパスでは
    デバイス自身がそれを行わねばならない．\margin{PCIe デバイスはこれを
    SR-IOV で行う．1枚のカードが多数の仮想ファンクションを見せ，それぞれが
    自前のレジスタとキューを持ち，別々のプロセスや仮想マシンに対応付け
    られる（第12章）．}
\end{itemize}

\section{同期アクセスと非同期アクセス}
\label{sec:l11-sync}

デバイスが要求を実行している間，\source{P3L5，同期アクセスと非同期アクセス；
Silberschatz ら ch.~12．}それを発行したスレッドは待つこともできるし，
実行を続けることもできる．

\begin{definition}[同期I/O]
  \term[どうきIO]{同期I/O}[synchronous I/O]では，I/O 要求を発行したスレッド
  は，その要求が完了するまでブロックされる．オペレーティングシステムは
  そのスレッドを，デバイスに対応する待ち行列，すなわち第2章の I/O キューに
  置き，デバイスが完了を報告したときに再び実行可能にする．
\end{definition}

\term*{非同期I/O}\index{ひどうきIO@非同期I/O}（第6章）では，呼び出しは直ちに
戻り，スレッドは実行を続ける（\cref{fig:l11-sync-async}）．スレッドは結果を
後から2つの方法のいずれかで知る．スレッド自身が操作の完了を確認して結果を
取り出すか，オペレーティングシステムが，操作が完了して結果が利用可能になった
ことを通知するかである．カーネルがデバイスへの非同期アクセスに対応していない
場合には，第6章の Flash ウェブサーバのように，主スレッドに代わってブロッキング
呼び出しを行うヘルパスレッドやヘルパプロセスを使って，ユーザ空間で模倣できる．%
\margin{glibc は POSIX の \texttt{aio\_read} をまさにこのヘルパスレッドに
よる模倣で実装している．\texttt{io\_uring}（第6章）は，それを不要にする
カーネルのインタフェースである．}

\begin{figure}[htbp]
  \centering
  % Synchronous versus asynchronous I/O: timelines of the requesting thread and
% the device.
% Usage: % Synchronous versus asynchronous I/O: timelines of the requesting thread and
% the device.
% Usage: % Synchronous versus asynchronous I/O: timelines of the requesting thread and
% the device.
% Usage: \input{figures/l11-sync-async}   (width ~116 mm)
\begin{tikzpicture}[x=1mm, y=1mm, font=\footnotesize\sffamily,
  >={Stealth[length=1.6mm]},
  lab/.style={font=\scriptsize\sffamily, align=center},
  row/.style={font=\scriptsize\sffamily, anchor=east},
  ttl/.style={font=\footnotesize\sffamily\bfseries, anchor=west}]
  % \lelseg{x1}{x2}{y}{fill}{text}
  \def\lelseg#1#2#3#4#5{%
    \draw[fill=#4] (#1,#3-2.5) rectangle (#2,#3+2.5);
    \node[lab] at ({(#1+#2)/2},#3) {#5};}
  % (a) synchronous
  \node[ttl] at (0,31) {\iflangja{(a) 同期I/O}{(a) synchronous I/O}};
  \node[row] at (20,24) {\iflangja{スレッド}{thread}};
  \lelseg{22}{44}{24}{white}{\iflangja{実行}{running}}
  \lelseg{44}{86}{24}{black!15}{\iflangja{ブロック（デバイスの待ち行列）}{blocked (on the device's wait queue)}}
  \lelseg{86}{116}{24}{white}{\iflangja{結果を使って続行}{continues with result}}
  \node[row] at (20,15) {\iflangja{デバイス}{device}};
  \lelseg{48}{80}{15}{ln-box}{\iflangja{I/O 処理}{I/O operation}}
  \draw[->] (44,21.5) -- node[lab, left=0.3mm, pos=0.35] {\iflangja{要求}{request}} (48,17.5);
  \draw[->, ln-caution] (80,17.5) -- node[lab, right=0.3mm, pos=0.55, text=ln-caution] {\iflangja{完了}{done}} (86,21.5);
  % (b) asynchronous
  \node[ttl] at (0,6) {\iflangja{(b) 非同期I/O}{(b) asynchronous I/O}};
  \node[row] at (20,-1) {\iflangja{スレッド}{thread}};
  \lelseg{22}{44}{-1}{white}{\iflangja{実行}{running}}
  \lelseg{44}{92}{-1}{white}{\iflangja{他の処理を続ける}{continues with other work}}
  \lelseg{92}{116}{-1}{ln-accent!22}{\iflangja{結果を取得}{retrieves result}}
  \node[row] at (20,-10) {\iflangja{デバイス}{device}};
  \lelseg{48}{80}{-10}{ln-box}{\iflangja{I/O 処理}{I/O operation}}
  \draw[->] (44,-3.5) -- node[lab, left=0.3mm, pos=0.35] {\iflangja{要求}{request}} (48,-7.5);
  \draw[->, densely dashed, ln-caution] (80,-7.5) -- (92,-3.5);
  \node[lab, anchor=north west, text=ln-caution] at (81,-12.8)
    {\iflangja{通知，またはスレッドが確認}{notification, or the thread checks}};
  % time axis
  \draw[->] (22,-21) -- node[lab, below] {\iflangja{時間}{time}} (116,-21);
\end{tikzpicture}
   (width ~116 mm)
\begin{tikzpicture}[x=1mm, y=1mm, font=\footnotesize\sffamily,
  >={Stealth[length=1.6mm]},
  lab/.style={font=\scriptsize\sffamily, align=center},
  row/.style={font=\scriptsize\sffamily, anchor=east},
  ttl/.style={font=\footnotesize\sffamily\bfseries, anchor=west}]
  % \lelseg{x1}{x2}{y}{fill}{text}
  \def\lelseg#1#2#3#4#5{%
    \draw[fill=#4] (#1,#3-2.5) rectangle (#2,#3+2.5);
    \node[lab] at ({(#1+#2)/2},#3) {#5};}
  % (a) synchronous
  \node[ttl] at (0,31) {\iflangja{(a) 同期I/O}{(a) synchronous I/O}};
  \node[row] at (20,24) {\iflangja{スレッド}{thread}};
  \lelseg{22}{44}{24}{white}{\iflangja{実行}{running}}
  \lelseg{44}{86}{24}{black!15}{\iflangja{ブロック（デバイスの待ち行列）}{blocked (on the device's wait queue)}}
  \lelseg{86}{116}{24}{white}{\iflangja{結果を使って続行}{continues with result}}
  \node[row] at (20,15) {\iflangja{デバイス}{device}};
  \lelseg{48}{80}{15}{ln-box}{\iflangja{I/O 処理}{I/O operation}}
  \draw[->] (44,21.5) -- node[lab, left=0.3mm, pos=0.35] {\iflangja{要求}{request}} (48,17.5);
  \draw[->, ln-caution] (80,17.5) -- node[lab, right=0.3mm, pos=0.55, text=ln-caution] {\iflangja{完了}{done}} (86,21.5);
  % (b) asynchronous
  \node[ttl] at (0,6) {\iflangja{(b) 非同期I/O}{(b) asynchronous I/O}};
  \node[row] at (20,-1) {\iflangja{スレッド}{thread}};
  \lelseg{22}{44}{-1}{white}{\iflangja{実行}{running}}
  \lelseg{44}{92}{-1}{white}{\iflangja{他の処理を続ける}{continues with other work}}
  \lelseg{92}{116}{-1}{ln-accent!22}{\iflangja{結果を取得}{retrieves result}}
  \node[row] at (20,-10) {\iflangja{デバイス}{device}};
  \lelseg{48}{80}{-10}{ln-box}{\iflangja{I/O 処理}{I/O operation}}
  \draw[->] (44,-3.5) -- node[lab, left=0.3mm, pos=0.35] {\iflangja{要求}{request}} (48,-7.5);
  \draw[->, densely dashed, ln-caution] (80,-7.5) -- (92,-3.5);
  \node[lab, anchor=north west, text=ln-caution] at (81,-12.8)
    {\iflangja{通知，またはスレッドが確認}{notification, or the thread checks}};
  % time axis
  \draw[->] (22,-21) -- node[lab, below] {\iflangja{時間}{time}} (116,-21);
\end{tikzpicture}
   (width ~116 mm)
\begin{tikzpicture}[x=1mm, y=1mm, font=\footnotesize\sffamily,
  >={Stealth[length=1.6mm]},
  lab/.style={font=\scriptsize\sffamily, align=center},
  row/.style={font=\scriptsize\sffamily, anchor=east},
  ttl/.style={font=\footnotesize\sffamily\bfseries, anchor=west}]
  % \lelseg{x1}{x2}{y}{fill}{text}
  \def\lelseg#1#2#3#4#5{%
    \draw[fill=#4] (#1,#3-2.5) rectangle (#2,#3+2.5);
    \node[lab] at ({(#1+#2)/2},#3) {#5};}
  % (a) synchronous
  \node[ttl] at (0,31) {\iflangja{(a) 同期I/O}{(a) synchronous I/O}};
  \node[row] at (20,24) {\iflangja{スレッド}{thread}};
  \lelseg{22}{44}{24}{white}{\iflangja{実行}{running}}
  \lelseg{44}{86}{24}{black!15}{\iflangja{ブロック（デバイスの待ち行列）}{blocked (on the device's wait queue)}}
  \lelseg{86}{116}{24}{white}{\iflangja{結果を使って続行}{continues with result}}
  \node[row] at (20,15) {\iflangja{デバイス}{device}};
  \lelseg{48}{80}{15}{ln-box}{\iflangja{I/O 処理}{I/O operation}}
  \draw[->] (44,21.5) -- node[lab, left=0.3mm, pos=0.35] {\iflangja{要求}{request}} (48,17.5);
  \draw[->, ln-caution] (80,17.5) -- node[lab, right=0.3mm, pos=0.55, text=ln-caution] {\iflangja{完了}{done}} (86,21.5);
  % (b) asynchronous
  \node[ttl] at (0,6) {\iflangja{(b) 非同期I/O}{(b) asynchronous I/O}};
  \node[row] at (20,-1) {\iflangja{スレッド}{thread}};
  \lelseg{22}{44}{-1}{white}{\iflangja{実行}{running}}
  \lelseg{44}{92}{-1}{white}{\iflangja{他の処理を続ける}{continues with other work}}
  \lelseg{92}{116}{-1}{ln-accent!22}{\iflangja{結果を取得}{retrieves result}}
  \node[row] at (20,-10) {\iflangja{デバイス}{device}};
  \lelseg{48}{80}{-10}{ln-box}{\iflangja{I/O 処理}{I/O operation}}
  \draw[->] (44,-3.5) -- node[lab, left=0.3mm, pos=0.35] {\iflangja{要求}{request}} (48,-7.5);
  \draw[->, densely dashed, ln-caution] (80,-7.5) -- (92,-3.5);
  \node[lab, anchor=north west, text=ln-caution] at (81,-12.8)
    {\iflangja{通知，またはスレッドが確認}{notification, or the thread checks}};
  % time axis
  \draw[->] (22,-21) -- node[lab, below] {\iflangja{時間}{time}} (116,-21);
\end{tikzpicture}

  \caption{(a) 同期I/O では，デバイスが操作を完了するまでスレッドはブロック
    する．(b) 非同期I/O では，スレッドは実行を続け，通知を受けた後または自ら
    確認した後に結果を取得する．}
  \label{fig:l11-sync-async}
\end{figure}

\begin{keypoint}
  デバイスアクセスは通常，システムコールからカーネル内スタックとドライバを
  経てデバイスに至り，結果は同じ経路を戻る．OS バイパスは，デバイスが自身の
  資源を保護し逆多重化できるとき，この経路からカーネルを取り除く．経路に
  かかわらず，要求したスレッドはブロックする（同期I/O）か，実行を続けて後から
  結果を回収する（非同期I/O）．
\end{keypoint}

\section{ブロックデバイススタック}
\label{sec:l11-block-stack}

プロセスがブロックデバイスをブロック番号で指定することはまれである．
\source{P3L5，ブロックデバイススタック；OSTEP ch.~36，39．}プロセスは，論理的
な記憶単位であるファイルを，\texttt{open}，\texttt{read}，\texttt{write}，
\texttt{close} といった操作を持つ POSIX API を通じて使う．API の下でカーネル
は層状に構成され，各層はその下の層のインタフェースを使う
（\cref{fig:l11-block-stack}）．

\begin{definition}[ファイルシステム]
  \term[ふぁいるしすてむ]{ファイルシステム}[file system]とは，記憶デバイスの
  ブロックをファイルとディレクトリへと編成するカーネルの構成要素である．
  各ファイルのデータがどこにあるか，またファイルをどう見つけてアクセスするか
  を決める．
\end{definition}

オペレーティングシステムは，すべてのファイルシステムが実装するインタフェース
を規定する．ファイルシステムの下には，ブロックデバイスどうしの差異を隠す層が
ある．

\begin{definition}[汎用ブロック層]
  \term[はんようぶろっくそう]{汎用ブロック層}[generic block layer]とは，
  任意のブロックデバイス上でブロックを読み書きするための，オペレーティング
  システムが標準化した単一のインタフェースを提供し，各要求を個々のデバイスの
  ドライバのインタフェースへ変換するカーネルの層である．
\end{definition}

したがって同じファイルシステムが，あらゆる種類のブロックデバイス上で動く．
汎用ブロック層が各要求をデバイスのドライバに合わせ，ドライバが SCSI や ATA
といったそのデバイスのプロトコルでハードウェアとやり取りする．

\begin{figure}[htbp]
  \centering
  % The block device stack: file system, generic block layer and device driver
% between a process and a block device.
% Usage: % The block device stack: file system, generic block layer and device driver
% between a process and a block device.
% Usage: % The block device stack: file system, generic block layer and device driver
% between a process and a block device.
% Usage: \input{figures/l11-block-stack}   (width ~112 mm)
\begin{tikzpicture}[x=1mm, y=1mm, font=\footnotesize\sffamily,
  >={Stealth[length=1.6mm]},
  box/.style={draw, fill=white, minimum width=38mm, minimum height=7mm, inner sep=1pt},
  kbox/.style={draw, fill=ln-box, minimum width=38mm, minimum height=7mm, inner sep=1pt},
  dev/.style={draw, fill=ln-accent!22, minimum width=38mm, minimum height=7mm, inner sep=1pt},
  ifc/.style={font=\scriptsize\sffamily, anchor=west, align=left},
  side/.style={font=\scriptsize\sffamily\itshape, text=ln-muted, anchor=west}]
  \node[box]  (p)  at (34,40) {\iflangja{ユーザプロセス}{user process}};
  \node[kbox] (fs) at (34,27) {\iflangja{ファイルシステム}{file system}};
  \node[kbox] (gb) at (34,14) {\iflangja{汎用ブロック層}{generic block layer}};
  \node[kbox] (dd) at (34,1)  {\iflangja{デバイスドライバ}{device driver}};
  \node[dev]  (bd) at (34,-12) {\iflangja{ブロックデバイス}{block device}};
  \foreach \a/\b in {p/fs, fs/gb, gb/dd, dd/bd} \draw[<->, thick] (\a.south) -- (\b.north);
  \node[ifc] at (40,33.5) {\iflangja{POSIX API：open, read, write, close}{POSIX API: open, read, write, close}};
  \node[ifc] at (40,20.5) {\iflangja{汎用ブロックインタフェース：ブロックの読み書き}{generic block interface: read or write blocks}};
  \node[ifc] at (40,7.5)  {\iflangja{プロトコル固有のインタフェース}{protocol-specific interface}};
  \node[ifc] at (40,-5.5) {\iflangja{デバイスのプロトコル（SCSI, ATA など）}{device protocol (e.g.\ SCSI, ATA)}};
  % boundaries
  \draw[densely dashed, ln-muted] (0,33.5) -- (15,33.5);
  \draw[densely dashed, ln-muted] (0,-5.5) -- (15,-5.5);
  \node[side, anchor=south west] at (0,33.8) {\iflangja{ユーザ}{user}};
  \node[side, anchor=north west] at (0,33.2) {\iflangja{カーネル}{kernel}};
  \node[side, anchor=north west] at (0,-5.8) {\iflangja{ハードウェア}{hardware}};
  % ioctl
  \draw[->, densely dashed, ln-accent] (p.west) .. controls (10,34) and (10,7) .. (dd.west);
  \node[font=\scriptsize\ttfamily, text=ln-accent, fill=white, inner sep=0.8pt] at (11.2,20.5) {ioctl};
\end{tikzpicture}
   (width ~112 mm)
\begin{tikzpicture}[x=1mm, y=1mm, font=\footnotesize\sffamily,
  >={Stealth[length=1.6mm]},
  box/.style={draw, fill=white, minimum width=38mm, minimum height=7mm, inner sep=1pt},
  kbox/.style={draw, fill=ln-box, minimum width=38mm, minimum height=7mm, inner sep=1pt},
  dev/.style={draw, fill=ln-accent!22, minimum width=38mm, minimum height=7mm, inner sep=1pt},
  ifc/.style={font=\scriptsize\sffamily, anchor=west, align=left},
  side/.style={font=\scriptsize\sffamily\itshape, text=ln-muted, anchor=west}]
  \node[box]  (p)  at (34,40) {\iflangja{ユーザプロセス}{user process}};
  \node[kbox] (fs) at (34,27) {\iflangja{ファイルシステム}{file system}};
  \node[kbox] (gb) at (34,14) {\iflangja{汎用ブロック層}{generic block layer}};
  \node[kbox] (dd) at (34,1)  {\iflangja{デバイスドライバ}{device driver}};
  \node[dev]  (bd) at (34,-12) {\iflangja{ブロックデバイス}{block device}};
  \foreach \a/\b in {p/fs, fs/gb, gb/dd, dd/bd} \draw[<->, thick] (\a.south) -- (\b.north);
  \node[ifc] at (40,33.5) {\iflangja{POSIX API：open, read, write, close}{POSIX API: open, read, write, close}};
  \node[ifc] at (40,20.5) {\iflangja{汎用ブロックインタフェース：ブロックの読み書き}{generic block interface: read or write blocks}};
  \node[ifc] at (40,7.5)  {\iflangja{プロトコル固有のインタフェース}{protocol-specific interface}};
  \node[ifc] at (40,-5.5) {\iflangja{デバイスのプロトコル（SCSI, ATA など）}{device protocol (e.g.\ SCSI, ATA)}};
  % boundaries
  \draw[densely dashed, ln-muted] (0,33.5) -- (15,33.5);
  \draw[densely dashed, ln-muted] (0,-5.5) -- (15,-5.5);
  \node[side, anchor=south west] at (0,33.8) {\iflangja{ユーザ}{user}};
  \node[side, anchor=north west] at (0,33.2) {\iflangja{カーネル}{kernel}};
  \node[side, anchor=north west] at (0,-5.8) {\iflangja{ハードウェア}{hardware}};
  % ioctl
  \draw[->, densely dashed, ln-accent] (p.west) .. controls (10,34) and (10,7) .. (dd.west);
  \node[font=\scriptsize\ttfamily, text=ln-accent, fill=white, inner sep=0.8pt] at (11.2,20.5) {ioctl};
\end{tikzpicture}
   (width ~112 mm)
\begin{tikzpicture}[x=1mm, y=1mm, font=\footnotesize\sffamily,
  >={Stealth[length=1.6mm]},
  box/.style={draw, fill=white, minimum width=38mm, minimum height=7mm, inner sep=1pt},
  kbox/.style={draw, fill=ln-box, minimum width=38mm, minimum height=7mm, inner sep=1pt},
  dev/.style={draw, fill=ln-accent!22, minimum width=38mm, minimum height=7mm, inner sep=1pt},
  ifc/.style={font=\scriptsize\sffamily, anchor=west, align=left},
  side/.style={font=\scriptsize\sffamily\itshape, text=ln-muted, anchor=west}]
  \node[box]  (p)  at (34,40) {\iflangja{ユーザプロセス}{user process}};
  \node[kbox] (fs) at (34,27) {\iflangja{ファイルシステム}{file system}};
  \node[kbox] (gb) at (34,14) {\iflangja{汎用ブロック層}{generic block layer}};
  \node[kbox] (dd) at (34,1)  {\iflangja{デバイスドライバ}{device driver}};
  \node[dev]  (bd) at (34,-12) {\iflangja{ブロックデバイス}{block device}};
  \foreach \a/\b in {p/fs, fs/gb, gb/dd, dd/bd} \draw[<->, thick] (\a.south) -- (\b.north);
  \node[ifc] at (40,33.5) {\iflangja{POSIX API：open, read, write, close}{POSIX API: open, read, write, close}};
  \node[ifc] at (40,20.5) {\iflangja{汎用ブロックインタフェース：ブロックの読み書き}{generic block interface: read or write blocks}};
  \node[ifc] at (40,7.5)  {\iflangja{プロトコル固有のインタフェース}{protocol-specific interface}};
  \node[ifc] at (40,-5.5) {\iflangja{デバイスのプロトコル（SCSI, ATA など）}{device protocol (e.g.\ SCSI, ATA)}};
  % boundaries
  \draw[densely dashed, ln-muted] (0,33.5) -- (15,33.5);
  \draw[densely dashed, ln-muted] (0,-5.5) -- (15,-5.5);
  \node[side, anchor=south west] at (0,33.8) {\iflangja{ユーザ}{user}};
  \node[side, anchor=north west] at (0,33.2) {\iflangja{カーネル}{kernel}};
  \node[side, anchor=north west] at (0,-5.8) {\iflangja{ハードウェア}{hardware}};
  % ioctl
  \draw[->, densely dashed, ln-accent] (p.west) .. controls (10,34) and (10,7) .. (dd.west);
  \node[font=\scriptsize\ttfamily, text=ln-accent, fill=white, inner sep=0.8pt] at (11.2,20.5) {ioctl};
\end{tikzpicture}

  \caption{ブロックデバイススタック．各層は，その下に示したインタフェースを
    使う．\texttt{ioctl} はデバイス固有の要求をドライバへ運ぶ．}
  \label{fig:l11-block-stack}
\end{figure}

デバイスに対する操作の中には，\texttt{read} にも \texttt{write} にも当て
はまらないものがある．Linux をはじめとする Unix 系システムは，デバイス固有の
情報にプロセスがアクセスし操作するためのシステムコール
\term[ioctl]{\texttt{ioctl}}[ioctl]（「I/O control」）を提供する．これは，
開いたデバイスファイルのファイルディスクリプタ（\cref{sec:l11-vfs-abs}），
デバイスのクラスまたはドライバに対して定義された要求コード，そして要求に
応じて要求固有のデータへのポインタを取る．

\needspace{11\baselineskip}
\begin{example}
  あるプログラムが1台目の光学ドライブのトレイを開く．
  \begin{lstlisting}[language=C, numbers=none]
#include <fcntl.h>
#include <unistd.h>
#include <sys/ioctl.h>
#include <linux/cdrom.h>

int fd = open("/dev/sr0", O_RDONLY | O_NONBLOCK);
if (fd >= 0) {
  ioctl(fd, CDROMEJECT);   /* read/write では表せない */
  close(fd);
}
\end{lstlisting}
  \texttt{O\_NONBLOCK} フラグにより，ディスクが入っていなくてもデバイス
  ファイルを開ける．
\end{example}

\section{仮想ファイルシステム}
\label{sec:l11-vfs}

ほとんどのシステムでは，単一のローカルデバイス上の単一のファイルシステム
では足りない．\source{P3L5，仮想ファイルシステム；OSTEP ch.~39；Kerrisk
ch.~14．}ファイルは複数のデバイスに存在しうるし，デバイスによって最も適した
ファイルシステムの実装も異なりうる．さらにファイルはローカルデバイス上に
まったく存在せず，ネットワーク越しにアクセスするリモートサーバ上にあるかも
しれない（第14章）．

\begin{definition}[仮想ファイルシステム]
  \term[かそうふぁいるしすてむ]{仮想ファイルシステム}[virtual file system]%
  \index{VFS|see{仮想ファイルシステム}}（VFS）とは，アプリケーションに1つの
  ファイルシステムインタフェースを見せ，各操作を，ファイルが存在する具体的な
  ファイルシステムへ対応付けるカーネルの層である．
\end{definition}

VFS は下層のファイルシステムの詳細をすべてアプリケーションから隠す
（\cref{fig:l11-vfs}）．アプリケーションはファイルがどこにあっても同じ API を
使い，VFS は各操作を該当するファイルシステム---ディスク上のローカルな ext2
ファイルシステムや，NFS のようなリモートファイルシステム---へ転送する．

\begin{figure}[tb]
  \centering
  % The virtual file system layer between the file system interface and the
% concrete local and remote file systems.
% Usage: % The virtual file system layer between the file system interface and the
% concrete local and remote file systems.
% Usage: % The virtual file system layer between the file system interface and the
% concrete local and remote file systems.
% Usage: \input{figures/l11-vfs}   (width ~114 mm)
\begin{tikzpicture}[x=1mm, y=1mm, font=\footnotesize\sffamily,
  >={Stealth[length=1.6mm]},
  bar/.style={draw, fill=ln-box, minimum width=108mm, minimum height=6.5mm, inner sep=1pt},
  fs/.style={draw, fill=white, minimum width=34mm, minimum height=9mm, inner sep=1pt, align=center},
  dev/.style={draw, fill=ln-accent!22, minimum width=22mm, minimum height=6mm, inner sep=1pt}]
  \node[draw, fill=white, minimum width=60mm, minimum height=6.5mm] (app) at (57,36)
    {\iflangja{ユーザプロセス}{user processes}};
  \node[bar] (fsi) at (57,25) {\iflangja{ファイルシステムインタフェース（POSIX API）}{file system interface (POSIX API)}};
  \node[bar, fill=ln-accent!10] (vfs) at (57,14) {\iflangja{仮想ファイルシステム（VFS）}{virtual file system (VFS)}};
  \node[fs] (f1) at (19,0) {\iflangja{ローカル FS\\（例：ext2）}{local file system\\(e.g.\ ext2)}};
  \node[fs] (f2) at (57,0) {\iflangja{ローカル FS\\（別の種類）}{local file system\\(another type)}};
  \node[fs] (f3) at (95,0) {\iflangja{リモート FS\\（例：NFS）}{remote file system\\(e.g.\ NFS)}};
  \node[dev] (d1) at (19,-14) {\iflangja{ディスク}{disk}};
  \node[dev] (d2) at (57,-14) {\iflangja{ディスク}{disk}};
  \node[dev] (d3) at (95,-14) {\iflangja{ネットワーク}{network}};
  \draw[<->, thick] (app) -- (fsi);
  \draw[<->, thick] (fsi) -- (vfs);
  \foreach \x/\f/\d in {19/f1/d1, 57/f2/d2, 95/f3/d3} {
    \draw[<->, thick] (\x,10.75) -- (\f.north);
    \draw[<->, thick] (\f.south) -- (\d.north);
  }
\end{tikzpicture}
   (width ~114 mm)
\begin{tikzpicture}[x=1mm, y=1mm, font=\footnotesize\sffamily,
  >={Stealth[length=1.6mm]},
  bar/.style={draw, fill=ln-box, minimum width=108mm, minimum height=6.5mm, inner sep=1pt},
  fs/.style={draw, fill=white, minimum width=34mm, minimum height=9mm, inner sep=1pt, align=center},
  dev/.style={draw, fill=ln-accent!22, minimum width=22mm, minimum height=6mm, inner sep=1pt}]
  \node[draw, fill=white, minimum width=60mm, minimum height=6.5mm] (app) at (57,36)
    {\iflangja{ユーザプロセス}{user processes}};
  \node[bar] (fsi) at (57,25) {\iflangja{ファイルシステムインタフェース（POSIX API）}{file system interface (POSIX API)}};
  \node[bar, fill=ln-accent!10] (vfs) at (57,14) {\iflangja{仮想ファイルシステム（VFS）}{virtual file system (VFS)}};
  \node[fs] (f1) at (19,0) {\iflangja{ローカル FS\\（例：ext2）}{local file system\\(e.g.\ ext2)}};
  \node[fs] (f2) at (57,0) {\iflangja{ローカル FS\\（別の種類）}{local file system\\(another type)}};
  \node[fs] (f3) at (95,0) {\iflangja{リモート FS\\（例：NFS）}{remote file system\\(e.g.\ NFS)}};
  \node[dev] (d1) at (19,-14) {\iflangja{ディスク}{disk}};
  \node[dev] (d2) at (57,-14) {\iflangja{ディスク}{disk}};
  \node[dev] (d3) at (95,-14) {\iflangja{ネットワーク}{network}};
  \draw[<->, thick] (app) -- (fsi);
  \draw[<->, thick] (fsi) -- (vfs);
  \foreach \x/\f/\d in {19/f1/d1, 57/f2/d2, 95/f3/d3} {
    \draw[<->, thick] (\x,10.75) -- (\f.north);
    \draw[<->, thick] (\f.south) -- (\d.north);
  }
\end{tikzpicture}
   (width ~114 mm)
\begin{tikzpicture}[x=1mm, y=1mm, font=\footnotesize\sffamily,
  >={Stealth[length=1.6mm]},
  bar/.style={draw, fill=ln-box, minimum width=108mm, minimum height=6.5mm, inner sep=1pt},
  fs/.style={draw, fill=white, minimum width=34mm, minimum height=9mm, inner sep=1pt, align=center},
  dev/.style={draw, fill=ln-accent!22, minimum width=22mm, minimum height=6mm, inner sep=1pt}]
  \node[draw, fill=white, minimum width=60mm, minimum height=6.5mm] (app) at (57,36)
    {\iflangja{ユーザプロセス}{user processes}};
  \node[bar] (fsi) at (57,25) {\iflangja{ファイルシステムインタフェース（POSIX API）}{file system interface (POSIX API)}};
  \node[bar, fill=ln-accent!10] (vfs) at (57,14) {\iflangja{仮想ファイルシステム（VFS）}{virtual file system (VFS)}};
  \node[fs] (f1) at (19,0) {\iflangja{ローカル FS\\（例：ext2）}{local file system\\(e.g.\ ext2)}};
  \node[fs] (f2) at (57,0) {\iflangja{ローカル FS\\（別の種類）}{local file system\\(another type)}};
  \node[fs] (f3) at (95,0) {\iflangja{リモート FS\\（例：NFS）}{remote file system\\(e.g.\ NFS)}};
  \node[dev] (d1) at (19,-14) {\iflangja{ディスク}{disk}};
  \node[dev] (d2) at (57,-14) {\iflangja{ディスク}{disk}};
  \node[dev] (d3) at (95,-14) {\iflangja{ネットワーク}{network}};
  \draw[<->, thick] (app) -- (fsi);
  \draw[<->, thick] (fsi) -- (vfs);
  \foreach \x/\f/\d in {19/f1/d1, 57/f2/d2, 95/f3/d3} {
    \draw[<->, thick] (\x,10.75) -- (\f.north);
    \draw[<->, thick] (\f.south) -- (\d.north);
  }
\end{tikzpicture}

  \caption{仮想ファイルシステムは，ファイルシステムインタフェースの操作を，
    具体的なローカルおよびリモートのファイルシステムへ対応付ける．}
  \label{fig:l11-vfs}
\end{figure}

\section{VFS の抽象化}
\label{sec:l11-vfs-abs}

具体的なファイルシステムはいずれも，VFS が定める少数の抽象化を実装する．
\source{P3L5，仮想ファイルシステムの抽象化，ディスク上の VFS；OSTEP
ch.~39--40；Kerrisk ch.~14．}VFS が操作する要素は\emph{ファイル}である．
プロセスは，開いたファイルをファイルディスクリプタによって参照する
（\cref{fig:l11-vfs-objects}）．

\begin{figure}[tb]
  \centering
  % VFS abstractions: file descriptor, file, dentries and inode in memory;
% superblock, inode blocks and data blocks on disk.
% Usage: % VFS abstractions: file descriptor, file, dentries and inode in memory;
% superblock, inode blocks and data blocks on disk.
% Usage: % VFS abstractions: file descriptor, file, dentries and inode in memory;
% superblock, inode blocks and data blocks on disk.
% Usage: \input{figures/l11-vfs-objects}   (width ~118 mm)
\begin{tikzpicture}[x=1mm, y=1mm, font=\scriptsize\sffamily,
  >={Stealth[length=1.5mm]},
  obj/.style={draw, fill=white, minimum height=8mm, inner sep=1pt, align=center},
  den/.style={draw, fill=white, minimum height=5mm, inner sep=1.2pt, font=\scriptsize\ttfamily},
  lab/.style={font=\scriptsize\sffamily, align=center},
  side/.style={font=\scriptsize\sffamily\bfseries, text=ln-accent}]
  \draw[densely dashed, ln-muted] (0,0) -- (118,0);
  \node[side, anchor=south east] at (118,0.4) {\iflangja{メモリ}{memory}};
  \node[side, anchor=north east] at (118,-0.4) {\iflangja{ディスク}{disk}};
  % file descriptor table of process P
  \node[lab, anchor=south] at (9,42.2) {\iflangja{P の fd}{fds of P}};
  \foreach \i/\y in {0/40,1/36,2/32,3/28} {
    \draw[fill=white] (4,\y-2) rectangle (14,\y+2);
    \node[font=\scriptsize\ttfamily] at (9,\y) {\i};
  }
  \draw[fill=ln-accent!22] (4,26) rectangle (14,30);
  \node[font=\scriptsize\ttfamily] at (9,28) {3};
  % file
  \node[obj, minimum width=26mm] (file) at (36,28)
    {\iflangja{ファイル\\（オフセット，モード）}{file\\(offset, mode)}};
  \draw[->] (14,28) -- (file.west);
  % dentries
  \draw[densely dotted, ln-muted, fill=ln-box] (3,9) rectangle (70,19);
  \node[lab, anchor=north west, text=ln-muted] at (3.5,8.6) {\iflangja{dentry キャッシュ（メモリ上のみ）}{dentry cache (in memory only)}};
  \node[den] (r) at (10,14) {/};
  \node[den] (u) at (24,14) {users};
  \node[den] (j) at (39,14) {joe};
  \node[den] (n) at (58,14) {notes.txt};
  \draw[->] (r) -- (u);  \draw[->] (u) -- (j);  \draw[->] (j) -- (n);
  \draw[->] (file.east) -| (n.north);
  % inode in memory
  \node[obj, minimum width=32mm, minimum height=15mm] (ino) at (97,20)
    {\iflangja{iノード\\所有者，権限，サイズ，\\ブロックの一覧}{inode\\owner, permissions, size,\\list of blocks}};
  \draw[->] (n.east) -- (ino.west |- n.east);
  % on-disk layout
  \def\seg#1#2#3#4{\draw[fill=#3] (#1,-15) rectangle (#2,-8);
    \node[lab] at ({(#1+#2)/2},-11.5) {#4};}
  \seg{4}{30}{ln-accent!22}{\iflangja{スーパーブロック}{superblock}}
  \seg{30}{58}{ln-box}{\iflangja{iノードブロック}{inode blocks}}
  \seg{58}{98}{white}{\iflangja{データブロック}{data blocks}}
  \seg{98}{118}{black!8}{\iflangja{空きブロック}{free blocks}}
  \draw[<->, densely dashed] (85,12.5) -- (48,-8);
  \node[lab, anchor=east] at (64,3.2) {\iflangja{ディスク上の表現}{on-disk copy}};
  \foreach \x in {66,78,90} \draw[->] (97,12.5) -- (\x,-8);
  \draw[->, ln-accent] (17,-8) .. controls (22,-4) and (36,-4) .. (44,-8);
  \node[lab, text=ln-accent, anchor=south] at (30,-3.7) {\iflangja{配置の地図}{map of the layout}};
\end{tikzpicture}
   (width ~118 mm)
\begin{tikzpicture}[x=1mm, y=1mm, font=\scriptsize\sffamily,
  >={Stealth[length=1.5mm]},
  obj/.style={draw, fill=white, minimum height=8mm, inner sep=1pt, align=center},
  den/.style={draw, fill=white, minimum height=5mm, inner sep=1.2pt, font=\scriptsize\ttfamily},
  lab/.style={font=\scriptsize\sffamily, align=center},
  side/.style={font=\scriptsize\sffamily\bfseries, text=ln-accent}]
  \draw[densely dashed, ln-muted] (0,0) -- (118,0);
  \node[side, anchor=south east] at (118,0.4) {\iflangja{メモリ}{memory}};
  \node[side, anchor=north east] at (118,-0.4) {\iflangja{ディスク}{disk}};
  % file descriptor table of process P
  \node[lab, anchor=south] at (9,42.2) {\iflangja{P の fd}{fds of P}};
  \foreach \i/\y in {0/40,1/36,2/32,3/28} {
    \draw[fill=white] (4,\y-2) rectangle (14,\y+2);
    \node[font=\scriptsize\ttfamily] at (9,\y) {\i};
  }
  \draw[fill=ln-accent!22] (4,26) rectangle (14,30);
  \node[font=\scriptsize\ttfamily] at (9,28) {3};
  % file
  \node[obj, minimum width=26mm] (file) at (36,28)
    {\iflangja{ファイル\\（オフセット，モード）}{file\\(offset, mode)}};
  \draw[->] (14,28) -- (file.west);
  % dentries
  \draw[densely dotted, ln-muted, fill=ln-box] (3,9) rectangle (70,19);
  \node[lab, anchor=north west, text=ln-muted] at (3.5,8.6) {\iflangja{dentry キャッシュ（メモリ上のみ）}{dentry cache (in memory only)}};
  \node[den] (r) at (10,14) {/};
  \node[den] (u) at (24,14) {users};
  \node[den] (j) at (39,14) {joe};
  \node[den] (n) at (58,14) {notes.txt};
  \draw[->] (r) -- (u);  \draw[->] (u) -- (j);  \draw[->] (j) -- (n);
  \draw[->] (file.east) -| (n.north);
  % inode in memory
  \node[obj, minimum width=32mm, minimum height=15mm] (ino) at (97,20)
    {\iflangja{iノード\\所有者，権限，サイズ，\\ブロックの一覧}{inode\\owner, permissions, size,\\list of blocks}};
  \draw[->] (n.east) -- (ino.west |- n.east);
  % on-disk layout
  \def\seg#1#2#3#4{\draw[fill=#3] (#1,-15) rectangle (#2,-8);
    \node[lab] at ({(#1+#2)/2},-11.5) {#4};}
  \seg{4}{30}{ln-accent!22}{\iflangja{スーパーブロック}{superblock}}
  \seg{30}{58}{ln-box}{\iflangja{iノードブロック}{inode blocks}}
  \seg{58}{98}{white}{\iflangja{データブロック}{data blocks}}
  \seg{98}{118}{black!8}{\iflangja{空きブロック}{free blocks}}
  \draw[<->, densely dashed] (85,12.5) -- (48,-8);
  \node[lab, anchor=east] at (64,3.2) {\iflangja{ディスク上の表現}{on-disk copy}};
  \foreach \x in {66,78,90} \draw[->] (97,12.5) -- (\x,-8);
  \draw[->, ln-accent] (17,-8) .. controls (22,-4) and (36,-4) .. (44,-8);
  \node[lab, text=ln-accent, anchor=south] at (30,-3.7) {\iflangja{配置の地図}{map of the layout}};
\end{tikzpicture}
   (width ~118 mm)
\begin{tikzpicture}[x=1mm, y=1mm, font=\scriptsize\sffamily,
  >={Stealth[length=1.5mm]},
  obj/.style={draw, fill=white, minimum height=8mm, inner sep=1pt, align=center},
  den/.style={draw, fill=white, minimum height=5mm, inner sep=1.2pt, font=\scriptsize\ttfamily},
  lab/.style={font=\scriptsize\sffamily, align=center},
  side/.style={font=\scriptsize\sffamily\bfseries, text=ln-accent}]
  \draw[densely dashed, ln-muted] (0,0) -- (118,0);
  \node[side, anchor=south east] at (118,0.4) {\iflangja{メモリ}{memory}};
  \node[side, anchor=north east] at (118,-0.4) {\iflangja{ディスク}{disk}};
  % file descriptor table of process P
  \node[lab, anchor=south] at (9,42.2) {\iflangja{P の fd}{fds of P}};
  \foreach \i/\y in {0/40,1/36,2/32,3/28} {
    \draw[fill=white] (4,\y-2) rectangle (14,\y+2);
    \node[font=\scriptsize\ttfamily] at (9,\y) {\i};
  }
  \draw[fill=ln-accent!22] (4,26) rectangle (14,30);
  \node[font=\scriptsize\ttfamily] at (9,28) {3};
  % file
  \node[obj, minimum width=26mm] (file) at (36,28)
    {\iflangja{ファイル\\（オフセット，モード）}{file\\(offset, mode)}};
  \draw[->] (14,28) -- (file.west);
  % dentries
  \draw[densely dotted, ln-muted, fill=ln-box] (3,9) rectangle (70,19);
  \node[lab, anchor=north west, text=ln-muted] at (3.5,8.6) {\iflangja{dentry キャッシュ（メモリ上のみ）}{dentry cache (in memory only)}};
  \node[den] (r) at (10,14) {/};
  \node[den] (u) at (24,14) {users};
  \node[den] (j) at (39,14) {joe};
  \node[den] (n) at (58,14) {notes.txt};
  \draw[->] (r) -- (u);  \draw[->] (u) -- (j);  \draw[->] (j) -- (n);
  \draw[->] (file.east) -| (n.north);
  % inode in memory
  \node[obj, minimum width=32mm, minimum height=15mm] (ino) at (97,20)
    {\iflangja{iノード\\所有者，権限，サイズ，\\ブロックの一覧}{inode\\owner, permissions, size,\\list of blocks}};
  \draw[->] (n.east) -- (ino.west |- n.east);
  % on-disk layout
  \def\seg#1#2#3#4{\draw[fill=#3] (#1,-15) rectangle (#2,-8);
    \node[lab] at ({(#1+#2)/2},-11.5) {#4};}
  \seg{4}{30}{ln-accent!22}{\iflangja{スーパーブロック}{superblock}}
  \seg{30}{58}{ln-box}{\iflangja{iノードブロック}{inode blocks}}
  \seg{58}{98}{white}{\iflangja{データブロック}{data blocks}}
  \seg{98}{118}{black!8}{\iflangja{空きブロック}{free blocks}}
  \draw[<->, densely dashed] (85,12.5) -- (48,-8);
  \node[lab, anchor=east] at (64,3.2) {\iflangja{ディスク上の表現}{on-disk copy}};
  \foreach \x in {66,78,90} \draw[->] (97,12.5) -- (\x,-8);
  \draw[->, ln-accent] (17,-8) .. controls (22,-4) and (36,-4) .. (44,-8);
  \node[lab, text=ln-accent, anchor=south] at (30,-3.7) {\iflangja{配置の地図}{map of the layout}};
\end{tikzpicture}

  \caption{VFS の抽象化．プロセス P のファイルディスクリプタ 3 は開いた
    ファイルを指し，そのファイルの dentry が iノードへ導く．パスの dentry は
    メモリ上にのみ存在する．ディスク上では，スーパーブロックが iノード・
    データ・空き領域を保持するブロックの地図となる．}
  \label{fig:l11-vfs-objects}
\end{figure}

\begin{definition}[ファイルディスクリプタ]
  \term[ふぁいるでぃすくりぷた]{ファイルディスクリプタ}[file descriptor]とは，
  プロセスにとっての開いたファイルの，オペレーティングシステムによる表現で
  ある．プロセスが \texttt{open} から得て，\texttt{read}，\texttt{write}，
  \texttt{sendfile}，ロック，\texttt{close} といったファイルへの操作に渡す
  小さな整数である．\margin{記述子 0，1，2 はプロセスの標準入力・標準出力・
  標準エラーである．既定では1つのプロセスが同時に開けるファイルは
  \num{1024} 個である（\texttt{ulimit -n}）．}
\end{definition}

ファイル自身の属性と，そのデータブロックの位置は，ファイルの iノードに記録
される．

\begin{definition}[iノード]
  \term[iのーど]{iノード}[inode]とは，ファイルの永続的な表現である．ファイル
  のすべてのデータブロックを列挙する索引であり，デバイス・権限・サイズと
  いったメタデータを伴う．%
  \caution{複数の名前が1つの iノードを共有でき（ハードリンク），名前の削除は
  リンク数を1減らすだけである．リンク数が0になり開いた記述子もなくなると
  ファイルは解放される．}
\end{definition}

ディレクトリもまたファイルであり，そのデータは，ディレクトリが含むファイルと
その iノードについての情報である．したがってファイル名はディレクトリに現れ，
iノードには現れない．パスの各構成要素は，さらに別のオブジェクトで表される．

\begin{definition}[ディレクトリエントリ]
  \term[でぃれくとりえんとり]{ディレクトリエントリ}[directory entry]%
  \index{dentry|see{ディレクトリエントリ}}（dentry）とは，パスの1つの構成
  要素に対応する VFS のオブジェクトである．
\end{definition}

パス \texttt{/users/joe} の解決では，\texttt{/}，\texttt{users}，\texttt{joe}
のそれぞれに対して dentry が作られる．VFS は dentry を \emph{dentry
キャッシュ}に保持するので，同じ接頭辞を持つパスを後で解決するときに，対応
するディレクトリを読み直さずに済む．dentry はソフトステートであり，メモリ上
にのみ存在してディスク上の表現を持たない．

\begin{definition}[スーパーブロック]
  \term[すーぱーぶろっく]{スーパーブロック}[superblock]は，ファイルシステム
  の配置に関する，ファイルシステム固有の情報を保持する．
\end{definition}

\begin{example}
  あるプロセスが \texttt{/users/joe/notes.txt} を開く．\texttt{/} の dentry
  から始めて，VFS はパス上の各ディレクトリを読んで次の構成要素とその iノード
  を見つけ，その過程で \texttt{users}，\texttt{joe}，\texttt{notes.txt} の
  dentry を作る．続いてそのプロセスが \texttt{/users/joe/todo.txt} を開くと，
  \texttt{/}，\texttt{users}，\texttt{joe} の dentry は dentry キャッシュで
  見つかり，最後の構成要素だけをディレクトリ \texttt{joe} の中で探せばよい．
\end{example}

\subsection{ディスク上の VFS}

ディスク上では，これらの抽象化は次のように配置される
（\cref{fig:l11-vfs-objects}）．ファイルのデータはデータブロックに保持される．
ファイルのブロックを管理する iノードもまた，ディスク上のいずれかのブロックに
置かれる．スーパーブロックは，どのブロックが iノードを保持し，どれがデータを
保持し，どれが空いているかという，ディスクブロック全体の地図を提供する．

\section{第2拡張ファイルシステム}
\label{sec:l11-ext2}

Linux は長らく，ディスク上の構造が VFS の抽象化と密接に対応するファイル
システムの一族を使ってきた．\source{P3L5，ext2；OSTEP ch.~40--41．}第2拡張
ファイルシステム \term{ext2}[ext2] は，長年にわたり Linux の既定のファイル
システムであり，その後継である ext3 と ext4 の基礎でもある．ext2 の
パーティションは，ブート用に予約された \SI{1024}{\byte}（ブロックが
\SI{1}{\kibi\byte} のときはブロック 0 全体）から始まり，残りはブロックの
グループに分割される（\cref{fig:l11-ext2}）．

\begin{definition}[ブロックグループ]
  \term[ぶろっくぐるーぷ]{ブロックグループ}[block group]とは，ext2 が
  ディスクパーティションを分割してできる，等しい大きさの連続したブロックの
  グループの1つである．各ブロックグループは，自身の iノードとデータブロック
  を，それらを記述するメタデータとともに保持する．
\end{definition}

各ブロックグループは次のものを含む．
\begin{description}
  \item[スーパーブロックのコピー] ファイルシステム全体の情報，すなわち
    iノードの数とディスクブロックの数，空き iノードと空きブロックの数，そして
    ブロックの大きさを記録する．
  \item[グループ記述子表のコピー] グループごとに1つの記述子を持ち，その
    グループのビットマップと iノード表の位置，および空きブロック数・空き
    iノード数・ディレクトリ数を記録する．
  \item[ブロックビットマップと iノードビットマップ] そのグループのどの
    ブロックとどの iノードが空いているかを管理する．
  \item[iノード表] ファイルシステムの iノードのうち連続した範囲を保持する．
    iノードは 1 から最大数まで番号付けされ，ファイルごとに1つある．
  \item[データブロック] ファイルのデータを保持する．
\end{description}
スーパーブロックとグループ記述子表はファイルシステム全体を記述する．通常
使われるのはブロックグループ 0 のコピーであり，他のグループのコピーは
バックアップとして働く．あるファイルの iノードとデータブロックを同じグループ
に置くことで，関連するブロックがディスク上で近くに配置され，ディスクヘッドが
それらの間を移動する距離が減る．\margin{グループはディスクの幾何学的な領域
ではない．1ブロック分のビットマップで管理できるだけのブロックを保持する．
ブロックが \SI{4}{\kibi\byte} なら $8 \times 4096 = \num{32768}$ ブロック，
すなわち \SI{128}{\mebi\byte} である．}

\begin{figure}[htbp]
  \centering
  % Layout of an ext2 partition: boot block and block groups; contents of one
% block group.
% Usage: % Layout of an ext2 partition: boot block and block groups; contents of one
% block group.
% Usage: % Layout of an ext2 partition: boot block and block groups; contents of one
% block group.
% Usage: \input{figures/l11-ext2}   (width ~116 mm)
\begin{tikzpicture}[x=1mm, y=1mm, font=\scriptsize\sffamily,
  lab/.style={font=\scriptsize\sffamily, align=center},
  note/.style={font=\scriptsize\sffamily, text=ln-muted, align=center,
               anchor=north, text width=17mm, inner sep=0.5pt}]
  % top strip
  \def\tseg#1#2#3#4{\draw[fill=#3] (#1,18) rectangle (#2,25);
    \node[lab] at ({(#1+#2)/2},21.5) {#4};}
  \tseg{2}{22}{black!8}{\iflangja{ブート}{boot}}
  \tseg{22}{50}{white}{\iflangja{ブロックグループ 0}{block group 0}}
  \tseg{50}{78}{ln-box}{\iflangja{ブロックグループ 1}{block group 1}}
  \tseg{78}{88}{white}{\ldots}
  \tseg{88}{116}{white}{\iflangja{ブロックグループ $N$}{block group $N$}}
  \node[lab, anchor=south] at (12,25.3) {\iflangja{ブロック 0}{block 0}};
  % zoom lines
  \draw[densely dashed, ln-muted] (50,18) -- (2,5);
  \draw[densely dashed, ln-muted] (78,18) -- (116,5);
  % block group detail
  \def\bseg#1#2#3#4{\draw[fill=#3] (#1,-4) rectangle (#2,5);
    \node[lab] at ({(#1+#2)/2},0.5) {#4};}
  \bseg{2}{21}{ln-accent!22}{\iflangja{スーパー\\ブロック}{superblock}}
  \bseg{21}{40}{ln-accent!10}{\iflangja{グループ\\記述子}{group\\descriptor}}
  \bseg{40}{58}{ln-box}{\iflangja{ブロック\\ビットマップ}{block\\bitmap}}
  \bseg{58}{76}{ln-box}{\iflangja{iノード\\ビットマップ}{inode\\bitmap}}
  \bseg{76}{94}{white}{\iflangja{iノード}{inodes}}
  \bseg{94}{116}{white}{\iflangja{データ\\ブロック}{data\\blocks}}
  \node[note] at (11.5,-5) {\iflangja{iノード数，ブロック数，空きブロックの開始位置}{no.\ of inodes and blocks, start of free blocks}};
  \node[note] at (30.5,-5) {\iflangja{ビットマップの位置，空き iノード数，ディレクトリ数}{bitmaps, free inodes, directories}};
  \node[note] at (49,-5) {\iflangja{空きブロックを記録}{tracks free blocks}};
  \node[note] at (67,-5) {\iflangja{空き iノードを記録}{tracks free inodes}};
  \node[note] at (85,-5) {\iflangja{1 から最大数まで，ファイルごとに 1 つ}{1 to max.\ number, one per file}};
  \node[note] at (105,-5) {\iflangja{ファイルのデータ}{file data}};
\end{tikzpicture}
   (width ~116 mm)
\begin{tikzpicture}[x=1mm, y=1mm, font=\scriptsize\sffamily,
  lab/.style={font=\scriptsize\sffamily, align=center},
  note/.style={font=\scriptsize\sffamily, text=ln-muted, align=center,
               anchor=north, text width=17mm, inner sep=0.5pt}]
  % top strip
  \def\tseg#1#2#3#4{\draw[fill=#3] (#1,18) rectangle (#2,25);
    \node[lab] at ({(#1+#2)/2},21.5) {#4};}
  \tseg{2}{22}{black!8}{\iflangja{ブート}{boot}}
  \tseg{22}{50}{white}{\iflangja{ブロックグループ 0}{block group 0}}
  \tseg{50}{78}{ln-box}{\iflangja{ブロックグループ 1}{block group 1}}
  \tseg{78}{88}{white}{\ldots}
  \tseg{88}{116}{white}{\iflangja{ブロックグループ $N$}{block group $N$}}
  \node[lab, anchor=south] at (12,25.3) {\iflangja{ブロック 0}{block 0}};
  % zoom lines
  \draw[densely dashed, ln-muted] (50,18) -- (2,5);
  \draw[densely dashed, ln-muted] (78,18) -- (116,5);
  % block group detail
  \def\bseg#1#2#3#4{\draw[fill=#3] (#1,-4) rectangle (#2,5);
    \node[lab] at ({(#1+#2)/2},0.5) {#4};}
  \bseg{2}{21}{ln-accent!22}{\iflangja{スーパー\\ブロック}{superblock}}
  \bseg{21}{40}{ln-accent!10}{\iflangja{グループ\\記述子}{group\\descriptor}}
  \bseg{40}{58}{ln-box}{\iflangja{ブロック\\ビットマップ}{block\\bitmap}}
  \bseg{58}{76}{ln-box}{\iflangja{iノード\\ビットマップ}{inode\\bitmap}}
  \bseg{76}{94}{white}{\iflangja{iノード}{inodes}}
  \bseg{94}{116}{white}{\iflangja{データ\\ブロック}{data\\blocks}}
  \node[note] at (11.5,-5) {\iflangja{iノード数，ブロック数，空きブロックの開始位置}{no.\ of inodes and blocks, start of free blocks}};
  \node[note] at (30.5,-5) {\iflangja{ビットマップの位置，空き iノード数，ディレクトリ数}{bitmaps, free inodes, directories}};
  \node[note] at (49,-5) {\iflangja{空きブロックを記録}{tracks free blocks}};
  \node[note] at (67,-5) {\iflangja{空き iノードを記録}{tracks free inodes}};
  \node[note] at (85,-5) {\iflangja{1 から最大数まで，ファイルごとに 1 つ}{1 to max.\ number, one per file}};
  \node[note] at (105,-5) {\iflangja{ファイルのデータ}{file data}};
\end{tikzpicture}
   (width ~116 mm)
\begin{tikzpicture}[x=1mm, y=1mm, font=\scriptsize\sffamily,
  lab/.style={font=\scriptsize\sffamily, align=center},
  note/.style={font=\scriptsize\sffamily, text=ln-muted, align=center,
               anchor=north, text width=17mm, inner sep=0.5pt}]
  % top strip
  \def\tseg#1#2#3#4{\draw[fill=#3] (#1,18) rectangle (#2,25);
    \node[lab] at ({(#1+#2)/2},21.5) {#4};}
  \tseg{2}{22}{black!8}{\iflangja{ブート}{boot}}
  \tseg{22}{50}{white}{\iflangja{ブロックグループ 0}{block group 0}}
  \tseg{50}{78}{ln-box}{\iflangja{ブロックグループ 1}{block group 1}}
  \tseg{78}{88}{white}{\ldots}
  \tseg{88}{116}{white}{\iflangja{ブロックグループ $N$}{block group $N$}}
  \node[lab, anchor=south] at (12,25.3) {\iflangja{ブロック 0}{block 0}};
  % zoom lines
  \draw[densely dashed, ln-muted] (50,18) -- (2,5);
  \draw[densely dashed, ln-muted] (78,18) -- (116,5);
  % block group detail
  \def\bseg#1#2#3#4{\draw[fill=#3] (#1,-4) rectangle (#2,5);
    \node[lab] at ({(#1+#2)/2},0.5) {#4};}
  \bseg{2}{21}{ln-accent!22}{\iflangja{スーパー\\ブロック}{superblock}}
  \bseg{21}{40}{ln-accent!10}{\iflangja{グループ\\記述子}{group\\descriptor}}
  \bseg{40}{58}{ln-box}{\iflangja{ブロック\\ビットマップ}{block\\bitmap}}
  \bseg{58}{76}{ln-box}{\iflangja{iノード\\ビットマップ}{inode\\bitmap}}
  \bseg{76}{94}{white}{\iflangja{iノード}{inodes}}
  \bseg{94}{116}{white}{\iflangja{データ\\ブロック}{data\\blocks}}
  \node[note] at (11.5,-5) {\iflangja{iノード数，ブロック数，空きブロックの開始位置}{no.\ of inodes and blocks, start of free blocks}};
  \node[note] at (30.5,-5) {\iflangja{ビットマップの位置，空き iノード数，ディレクトリ数}{bitmaps, free inodes, directories}};
  \node[note] at (49,-5) {\iflangja{空きブロックを記録}{tracks free blocks}};
  \node[note] at (67,-5) {\iflangja{空き iノードを記録}{tracks free inodes}};
  \node[note] at (85,-5) {\iflangja{1 から最大数まで，ファイルごとに 1 つ}{1 to max.\ number, one per file}};
  \node[note] at (105,-5) {\iflangja{ファイルのデータ}{file data}};
\end{tikzpicture}

  \caption{ext2 のパーティションの配置．ブート用に予約された領域の後，
    パーティションはブロックグループに分割され，各グループはスーパーブロック
    とグループ記述子表のコピー，自身のビットマップ，そして iノードとデータ
    ブロックを持つ．}
  \label{fig:l11-ext2}
\end{figure}

\needspace{6\baselineskip}
\section{iノード}
\label{sec:l11-inodes}

ファイルは iノードで識別され，iノードはファイルのすべてのディスクブロックの
索引である．ブロックを順に列挙し，メタデータを伴う．\source{P3L5，iノード，
間接ポインタを持つ iノード；OSTEP ch.~40；Kerrisk ch.~14．}最も単純な形では，
iノードはファイルの全ブロックのブロック番号を保持する．このとき逐次アクセス
もランダムアクセスも同じくらい容易である．ブロックサイズを $B$ とすると，
オフセット $o$ のバイトは，一覧のエントリ $\lfloor o/B \rfloor$ が示す番号
のブロックにある．しかしファイルの大きさは，iノードに収まるエントリ
の数で制限される．

\begin{example}
  512 バイトの iノードがメタデータに 64 バイトを使うとすると，4 バイトの
  ブロック番号を $448/4 = 112$ 個収められる．\SI{4}{\kibi\byte} のブロック
  では，どのファイルも $112 \times \SI{4}{\kibi\byte} = \SI{448}{\kibi\byte}$
  より大きくなれない．
\end{example}

\subsection{間接ポインタ}

小さな iノードで大きなファイルを扱うために，iノードは数種類のポインタを持つ
（\cref{fig:l11-inode}）．
\begin{description}
  \item[直接ポインタ] \term[ちょくせつぽいんた]{直接ポインタ}[direct pointer]
    は，データブロックを指す．
  \item[間接ポインタ] \term[かんせつぽいんた]{間接ポインタ}[indirect pointer]
    は，データブロックへのポインタを含むブロックを指す．
  \item[二重間接ポインタ]
    \term[にじゅうかんせつぽいんた]{二重間接ポインタ}[double indirect pointer]
    は，データブロックへのポインタのブロックを指すポインタの，そのブロックを
    指す．
\end{description}
\emph{三重間接ポインタ}はさらにもう1段を加える．ext2 と多くの Unix 系
ファイルシステムの iノードは，12 個の直接ポインタと，各種類の間接ポインタを
1つずつ含む．\margin{ext2 の実際の値は，iノードが \SI{128}{\byte} で
パーティション \SI{16}{\kibi\byte} ごとに1つ，ブロック番号が4バイトであり，
ブロックが \SI{4}{\kibi\byte} なら $P = 1024$ となる
（\texttt{mke2fs.conf(5)}）．}

\begin{figure}[htbp]
  \centering
  % An inode with direct, single, double and triple indirect pointers.
% Usage: % An inode with direct, single, double and triple indirect pointers.
% Usage: % An inode with direct, single, double and triple indirect pointers.
% Usage: \input{figures/l11-inode}   (width ~116 mm)
\begin{tikzpicture}[x=1mm, y=1mm, font=\scriptsize\sffamily,
  >={Stealth[length=1.4mm]},
  data/.style={draw, fill=white, minimum width=10mm, minimum height=4.5mm,
               inner sep=0.5pt, font=\scriptsize\sffamily},
  lab/.style={font=\scriptsize\sffamily, align=center}]
  % \ptrblk{name}{x}{y}: block of pointers
  \def\ptrblk#1#2#3{%
    \draw[fill=ln-box] (#2-5,#3-4) rectangle (#2+5,#3+4);
    \foreach \k in {-2,0,2} \draw[ln-rule] (#2-5,#3+\k) -- (#2+5,#3+\k);
    \coordinate (#1w) at (#2-5,#3);
    \coordinate (#1e) at (#2+5,#3+3);
    \coordinate (#1s) at (#2+5,#3-3);}
  % inode
  \node[font=\footnotesize\sffamily\bfseries, anchor=south] at (19,48.3) {\iflangja{iノード}{inode}};
  \draw[fill=ln-accent!10] (4,36) rectangle (34,48);
  \node[lab] at (19,42) {\iflangja{メタデータ\\（モード，所有者，\\タイムスタンプ，サイズ）}{metadata\\(mode, owners,\\timestamps, size, \ldots)}};
  \def\irow#1#2{\draw[fill=white] (4,#1-2.5) rectangle (34,#1+2.5);
    \node[lab] at (19,#1) {#2};}
  \irow{33.5}{\iflangja{直接ポインタ 0}{direct 0}}
  \irow{28.5}{\ldots}
  \irow{23.5}{\iflangja{直接ポインタ 11}{direct 11}}
  \irow{18.5}{\iflangja{間接ポインタ}{single indirect}}
  \irow{13.5}{\iflangja{二重間接ポインタ}{double indirect}}
  \irow{8.5}{\iflangja{三重間接ポインタ}{triple indirect}}
  % direct
  \node[data] (a) at (111,33.5) {\iflangja{データ}{data}};
  \node[data] (b) at (111,23.5) {\iflangja{データ}{data}};
  \draw[->] (31,33.5) -- (a.west);
  \draw[->] (31,23.5) -- (b.west);
  % single indirect
  \ptrblk{p}{53}{15}
  \node[data] (c) at (111,15) {\iflangja{データ}{data}};
  \draw[->] (31,18.5) -- (pw);
  \draw[->] (pe) -- (c.west);
  % double indirect
  \ptrblk{q}{53}{0}
  \ptrblk{r}{77}{0}
  \node[data] (d) at (111,0) {\iflangja{データ}{data}};
  \draw[->] (31,13.5) -- (qw);
  \draw[->] (qe) -- (rw);
  \draw[->] (re) -- (d.west);
  % triple indirect
  \ptrblk{s}{53}{-15}
  \ptrblk{t}{71}{-15}
  \ptrblk{u}{89}{-15}
  \node[data] (e) at (111,-15) {\iflangja{データ}{data}};
  \draw[->] (31,8.5) -- (sw);
  \draw[->] (se) -- (tw);
  \draw[->] (te) -- (uw);
  \draw[->] (ue) -- (e.west);
  % legend
  \draw[fill=ln-box] (4,-27) rectangle (10,-23);
  \node[lab, anchor=west] at (11,-25) {\iflangja{ポインタのブロック}{block of pointers}};
  \draw[fill=white] (44,-27) rectangle (50,-23);
  \node[lab, anchor=west] at (51,-25) {\iflangja{データブロック}{data block}};
\end{tikzpicture}
   (width ~116 mm)
\begin{tikzpicture}[x=1mm, y=1mm, font=\scriptsize\sffamily,
  >={Stealth[length=1.4mm]},
  data/.style={draw, fill=white, minimum width=10mm, minimum height=4.5mm,
               inner sep=0.5pt, font=\scriptsize\sffamily},
  lab/.style={font=\scriptsize\sffamily, align=center}]
  % \ptrblk{name}{x}{y}: block of pointers
  \def\ptrblk#1#2#3{%
    \draw[fill=ln-box] (#2-5,#3-4) rectangle (#2+5,#3+4);
    \foreach \k in {-2,0,2} \draw[ln-rule] (#2-5,#3+\k) -- (#2+5,#3+\k);
    \coordinate (#1w) at (#2-5,#3);
    \coordinate (#1e) at (#2+5,#3+3);
    \coordinate (#1s) at (#2+5,#3-3);}
  % inode
  \node[font=\footnotesize\sffamily\bfseries, anchor=south] at (19,48.3) {\iflangja{iノード}{inode}};
  \draw[fill=ln-accent!10] (4,36) rectangle (34,48);
  \node[lab] at (19,42) {\iflangja{メタデータ\\（モード，所有者，\\タイムスタンプ，サイズ）}{metadata\\(mode, owners,\\timestamps, size, \ldots)}};
  \def\irow#1#2{\draw[fill=white] (4,#1-2.5) rectangle (34,#1+2.5);
    \node[lab] at (19,#1) {#2};}
  \irow{33.5}{\iflangja{直接ポインタ 0}{direct 0}}
  \irow{28.5}{\ldots}
  \irow{23.5}{\iflangja{直接ポインタ 11}{direct 11}}
  \irow{18.5}{\iflangja{間接ポインタ}{single indirect}}
  \irow{13.5}{\iflangja{二重間接ポインタ}{double indirect}}
  \irow{8.5}{\iflangja{三重間接ポインタ}{triple indirect}}
  % direct
  \node[data] (a) at (111,33.5) {\iflangja{データ}{data}};
  \node[data] (b) at (111,23.5) {\iflangja{データ}{data}};
  \draw[->] (31,33.5) -- (a.west);
  \draw[->] (31,23.5) -- (b.west);
  % single indirect
  \ptrblk{p}{53}{15}
  \node[data] (c) at (111,15) {\iflangja{データ}{data}};
  \draw[->] (31,18.5) -- (pw);
  \draw[->] (pe) -- (c.west);
  % double indirect
  \ptrblk{q}{53}{0}
  \ptrblk{r}{77}{0}
  \node[data] (d) at (111,0) {\iflangja{データ}{data}};
  \draw[->] (31,13.5) -- (qw);
  \draw[->] (qe) -- (rw);
  \draw[->] (re) -- (d.west);
  % triple indirect
  \ptrblk{s}{53}{-15}
  \ptrblk{t}{71}{-15}
  \ptrblk{u}{89}{-15}
  \node[data] (e) at (111,-15) {\iflangja{データ}{data}};
  \draw[->] (31,8.5) -- (sw);
  \draw[->] (se) -- (tw);
  \draw[->] (te) -- (uw);
  \draw[->] (ue) -- (e.west);
  % legend
  \draw[fill=ln-box] (4,-27) rectangle (10,-23);
  \node[lab, anchor=west] at (11,-25) {\iflangja{ポインタのブロック}{block of pointers}};
  \draw[fill=white] (44,-27) rectangle (50,-23);
  \node[lab, anchor=west] at (51,-25) {\iflangja{データブロック}{data block}};
\end{tikzpicture}
   (width ~116 mm)
\begin{tikzpicture}[x=1mm, y=1mm, font=\scriptsize\sffamily,
  >={Stealth[length=1.4mm]},
  data/.style={draw, fill=white, minimum width=10mm, minimum height=4.5mm,
               inner sep=0.5pt, font=\scriptsize\sffamily},
  lab/.style={font=\scriptsize\sffamily, align=center}]
  % \ptrblk{name}{x}{y}: block of pointers
  \def\ptrblk#1#2#3{%
    \draw[fill=ln-box] (#2-5,#3-4) rectangle (#2+5,#3+4);
    \foreach \k in {-2,0,2} \draw[ln-rule] (#2-5,#3+\k) -- (#2+5,#3+\k);
    \coordinate (#1w) at (#2-5,#3);
    \coordinate (#1e) at (#2+5,#3+3);
    \coordinate (#1s) at (#2+5,#3-3);}
  % inode
  \node[font=\footnotesize\sffamily\bfseries, anchor=south] at (19,48.3) {\iflangja{iノード}{inode}};
  \draw[fill=ln-accent!10] (4,36) rectangle (34,48);
  \node[lab] at (19,42) {\iflangja{メタデータ\\（モード，所有者，\\タイムスタンプ，サイズ）}{metadata\\(mode, owners,\\timestamps, size, \ldots)}};
  \def\irow#1#2{\draw[fill=white] (4,#1-2.5) rectangle (34,#1+2.5);
    \node[lab] at (19,#1) {#2};}
  \irow{33.5}{\iflangja{直接ポインタ 0}{direct 0}}
  \irow{28.5}{\ldots}
  \irow{23.5}{\iflangja{直接ポインタ 11}{direct 11}}
  \irow{18.5}{\iflangja{間接ポインタ}{single indirect}}
  \irow{13.5}{\iflangja{二重間接ポインタ}{double indirect}}
  \irow{8.5}{\iflangja{三重間接ポインタ}{triple indirect}}
  % direct
  \node[data] (a) at (111,33.5) {\iflangja{データ}{data}};
  \node[data] (b) at (111,23.5) {\iflangja{データ}{data}};
  \draw[->] (31,33.5) -- (a.west);
  \draw[->] (31,23.5) -- (b.west);
  % single indirect
  \ptrblk{p}{53}{15}
  \node[data] (c) at (111,15) {\iflangja{データ}{data}};
  \draw[->] (31,18.5) -- (pw);
  \draw[->] (pe) -- (c.west);
  % double indirect
  \ptrblk{q}{53}{0}
  \ptrblk{r}{77}{0}
  \node[data] (d) at (111,0) {\iflangja{データ}{data}};
  \draw[->] (31,13.5) -- (qw);
  \draw[->] (qe) -- (rw);
  \draw[->] (re) -- (d.west);
  % triple indirect
  \ptrblk{s}{53}{-15}
  \ptrblk{t}{71}{-15}
  \ptrblk{u}{89}{-15}
  \node[data] (e) at (111,-15) {\iflangja{データ}{data}};
  \draw[->] (31,8.5) -- (sw);
  \draw[->] (se) -- (tw);
  \draw[->] (te) -- (uw);
  \draw[->] (ue) -- (e.west);
  % legend
  \draw[fill=ln-box] (4,-27) rectangle (10,-23);
  \node[lab, anchor=west] at (11,-25) {\iflangja{ポインタのブロック}{block of pointers}};
  \draw[fill=white] (44,-27) rectangle (50,-23);
  \node[lab, anchor=west] at (51,-25) {\iflangja{データブロック}{data block}};
\end{tikzpicture}

  \caption{12 個の直接ポインタと，単一・二重・三重の間接ポインタを持つ
    iノード．間接の段が1つ増えるごとに，到達できるデータブロックの数は
    1ブロックあたりのポインタ数倍になる．}
  \label{fig:l11-inode}
\end{figure}

ブロックサイズを $B$，ポインタの大きさを $p$ とすると，1ブロックには
$P = B/p$ 個のポインタが入る．$n_d$ 個の直接ポインタと，3種類の間接ポインタ
を1つずつ持つ iノードが扱えるファイルの大きさは，最大で
\begin{equation}
  S_{\max} \;=\; \bigl(n_d + P + P^2 + P^3\bigr) \times B .
  \label{eq:l11-max-size}
\end{equation}

\begin{example}
  \SI{4}{\kibi\byte} のブロックと 8 バイトのポインタでは，1ブロックに
  $P = 4096/8 = 512$ 個のポインタが入る．12 個の直接ポインタは
  $12 \times \SI{4}{\kibi\byte} = \SI{48}{\kibi\byte}$ を，間接ポインタは
  $512 \times \SI{4}{\kibi\byte} = \SI{2}{\mebi\byte}$ を，二重間接ポインタは
  $512^2 \times \SI{4}{\kibi\byte} = \SI{1}{\gibi\byte}$ を，三重間接
  ポインタは $512^3 \times \SI{4}{\kibi\byte} = \SI{512}{\gibi\byte}$ を
  扱える．最大ファイルサイズはおよそ \SI{513}{\gibi\byte} である．
\end{example}

間接ポインタの利点は，小さな iノードで非常に大きなファイルを扱えることで
ある．代償はファイルアクセスが遅くなることであり，間接の段が1つ増えるごとに，
データブロックに到達するまでのディスクアクセスが1回増える．直接ポインタを
通じてブロックを読むには，iノードに1回とデータブロックに1回の計2回の
ディスクアクセスが必要であり，二重間接ポインタを通じて読むには最大4回が必要
である．\tip{ファイルのブロック $i$ にどのポインタで到達するかは，$i$ を
累積和 $n_d$，$n_d + P$，$n_d + P + P^2$ と比べれば分かる．最初に下回った和が
ポインタの種類を示す．}

\begin{example}
  前の例のパラメータで，あるプロセスがファイルのオフセット
  \SI{600}{\mebi\byte} のバイトを読む．これはファイルの
  ブロック $600 \times 1024 / 4 = \num{153600}$ にある．ブロック 0--11
  には直接に，ブロック 12--523 には間接ポインタを通じて，ブロック
  \num{524}--\num{262667} には二重間接ポインタを通じて到達するので，この
  読み出しには iノードと，2つのポインタのブロックと，データブロックが必要で
  あり，いずれもキャッシュされていなければ4回のディスクアクセスとなる．
\end{example}

\begin{keypoint}
  VFS は，ファイルディスクリプタ，dentry，iノード，スーパーブロックを通じて
  ファイルを操作する．ディスク上でのファイルは，その iノードと，iノードが
  索引付けするデータブロックである．直接ポインタと間接ポインタにより，固定長
  の iノードで非常に大きなファイルを索引付けできるが，間接の段ごとに追加の
  ディスクアクセスという代償を払う．
\end{keypoint}

\section{ディスクアクセスの最適化}
\label{sec:l11-optim}

ファイルへのアクセスにはディスクアクセスのコストがかかり，ハードディスクでは
逐次的でないアクセスのたびにヘッドを動かす時間がかかる．\source{P3L5，
ディスクアクセスの最適化；OSTEP ch.~37，40，42；Kerrisk ch.~13--14．}
ファイルシステムは，ファイルアクセスのオーバーヘッドを減らすために4つの技法
を用いる（\cref{tab:l11-optim}）．\margin{これらの技法が埋めようとする差は
大きい．\num{7200}~rpm のディスクは逐次読み出しでは
\SI{100}{\mega\byte\per\second} を超えるが，\SI{4}{\kibi\byte} の
ランダムアクセスは毎秒 100 回程度で，毎秒数百キロバイトにしかならない．}

\begin{table}[htbp]
  \centering
  \caption{ファイルアクセスのオーバーヘッドを減らす技法．}
  \label{tab:l11-optim}
  \small
  \begin{tabular}{@{}l>{\raggedright\arraybackslash}p{0.56\linewidth}@{}}
    \toprule
    技法 & 効果 \\
    \midrule
    キャッシングとバッファリング & ディスクアクセスの削減 \\
    I/O スケジューリング & ヘッド移動の削減．逐次アクセスが増え，ランダム
      アクセスが減る \\
    プリフェッチ & 局所性を利用してキャッシュヒットを増やす \\
    ジャーナリング & 更新をメモリだけに留めることなく，書き込みのランダム
      アクセスを減らす \\
    \bottomrule
  \end{tabular}
\end{table}

\subsection{キャッシングとバッファリング}

\begin{definition}[バッファキャッシュ]
  \term[ばっふぁきゃっしゅ]{バッファキャッシュ}[buffer cache]とは，主メモリ
  上のファイルデータのキャッシュである．読み出しは可能な限りキャッシュから
  提供され，書き込みはキャッシュを更新する．変更されたデータは定期的に，
  あるいはプロセスが \texttt{fsync} で要求したときに，ディスクへ書き出される．%
  \margin{このキャッシュの大きさは固定ではない．Linux ではページキャッシュ
  であり，第8章の置換方針が残したメモリをそのまま使う．}
\end{definition}

更新を確実に永続化しなければならないプロセスは，書き込みに続けて
\texttt{fsync} を呼ぶ．
\needspace{5\baselineskip}
\begin{lstlisting}[language=C, numbers=none]
write(fd, rec, sizeof *rec);  /* キャッシュを更新 */
fsync(fd);                    /* ディスクへ書き出す */
\end{lstlisting}
\caution{\texttt{write} が成功しても，データがディスク上にあるとは限らない．
クラッシュすると，まだキャッシュにあるデータは失われる．}
\texttt{fsync} は，ファイルのデータがデバイスへ書き込まれた後にのみ戻る．

\subsection{I/O スケジューリング}

ディスクヘッドの移動を減らすように，保留中のディスク要求を並べ替えることを
\term[IOすけじゅーりんぐ]{I/Oスケジューリング}[I/O scheduling]と呼ぶ．要求は，
逐次アクセスを最大にしランダムアクセスを最小にする順序で処理される．

\begin{example}
  ヘッドがブロック 10 に位置しており，ブロック 90，12，91，13，50 への
  書き込みがこの順に到着する．ブロック番号がディスク上の位置を反映すると
  仮定すると，到着順に処理した場合ヘッドは $80+78+79+78+37 = 352$ ブロック
  分を移動する．12，13，50，90，91 の順に処理すれば移動は
  $2+1+37+40+1 = 81$ ブロック分で済み，12 と 13，90 と 91 の組は逐次
  書き込みになる．\margin{ブロック番号順への並べ替えはエレベータ
  （SCAN）アルゴリズムである．Linux は \texttt{mq-deadline}，BFQ，Kyber を
  備え，NVMe では通常 \texttt{none} を使う．}
\end{example}

\needspace{5\baselineskip}
\subsection{プリフェッチ}

要求されたよりも多くのブロックを読むことを
\term[ぷりふぇっち]{プリフェッチ}[prefetching]あるいは先読みと呼ぶ．
あるブロックを読むときに，次に必要になりそうなブロックを一緒に読む．
プリフェッチは局所性---逐次的に読まれているファイルは間もなく後続のブロック
を必要とする---を利用し，バッファキャッシュのヒット率を高める．プロセスが
ファイルのブロック 40 を読み，ファイルシステムが同じディスク操作でブロック
41--43 も読めば，続く3回の読み出しはヘッドを動かさずにメモリから提供される．%
\margin{Linux の既定の先読み量は \SI{128}{\kibi\byte} であり，ブロック
デバイスごとに \texttt{/sys/block/*/queue/read\_ahead\_kb} で設定される．}

\subsection{ジャーナリング}

すべての更新を直ちに本来の位置へ書くとランダムアクセスが生じ，更新を
バッファキャッシュだけに留めると失う危険がある．
\term[じゃーなりんぐ]{ジャーナリング}[journaling]，すなわちロギングの技法は，
この両方の問題を避ける．ファイルシステムは各更新---どのブロックの，どの
オフセットに，どの値を書くか---を\emph{ジャーナル}と呼ばれるログに記述し，
ディスクの専用領域へ逐次的に書き，記録した更新を定期的に本来のディスク上の
位置へ反映する（\cref{fig:l11-journal}）．ジャーナルはディスク上にあるので，
そこに記録された更新はクラッシュを生き延び，後から反映できる．%
\caution{ジャーナリングは書き込みの総量を減らすわけではない．記録したブロック
は2回書かれる．散在した書き込みを，1回の逐次書き込みと後でまとめて行う更新に
置き換えるのである．}ext3 と ext4 のファイルシステムは，ext2 の設計にこのような
ジャーナルを加えたものである．
既定ではメタデータブロックだけをログに記録し，ファイルデータは最終的な位置へ
直接書く．データもログに記録することは選択肢として用意されている．

\begin{figure}[htbp]
  \centering
  % Journaling: updates are appended sequentially to the journal and later
% applied to their proper locations on disk.
% Usage: % Journaling: updates are appended sequentially to the journal and later
% applied to their proper locations on disk.
% Usage: % Journaling: updates are appended sequentially to the journal and later
% applied to their proper locations on disk.
% Usage: \input{figures/l11-journal}   (width ~116 mm)
\begin{tikzpicture}[x=1mm, y=1mm, font=\scriptsize\sffamily,
  >={Stealth[length=1.5mm]},
  lab/.style={font=\scriptsize\sffamily, align=center},
  num/.style={circle, draw=black, fill=white, text=black, line width=0.4pt,
              inner sep=0.4pt, font=\scriptsize\sffamily\bfseries}]
  % updates
  \node[draw, fill=white, align=center, inner sep=1.5pt] (up) at (39,22)
    {\iflangja{更新：ブロック 870，12，301}{updates to blocks 870, 12, 301}};
  \draw[->, line width=1.2pt, ln-accent] (up.south) -- node[num, right=0.8mm] {1} (39,8.2);
  \node[lab, anchor=west] at (43,14) {\iflangja{順に追記（シーケンシャル書き込み）}{appended in order (sequential writes)}};
  % disk strip
  \node[lab, anchor=east] at (17,4) {\iflangja{ディスク}{disk}};
  \draw[fill=white] (18,0) rectangle (116,8);
  \foreach \x/\b in {18/870, 32/12, 46/301} {
    \draw[fill=ln-accent!10] (\x,0) rectangle (\x+14,8);
    \node[font=\scriptsize\ttfamily] at (\x+7,4) {\b};
  }
  \foreach \x/\b in {66/12, 84/301, 102/870} {
    \draw[fill=ln-box] (\x,0) rectangle (\x+10,8);
    \node[font=\scriptsize\ttfamily] at (\x+5,4) {\b};
  }
  \node[lab, anchor=south west, text=ln-muted] at (18,8.3) {\iflangja{ジャーナル}{journal}};
  \node[lab, anchor=south, text=ln-muted] at (89,8.3) {\iflangja{本来の位置}{proper locations}};
  % apply later
  \draw[->, densely dashed] (25,0) .. controls (30,-15) and (100,-15) .. (107,0);
  \draw[->, densely dashed] (39,0) .. controls (42,-7) and (68,-7) .. (71,0);
  \draw[->, densely dashed] (53,0) .. controls (58,-10) and (84,-10) .. (89,0);
  \node[num] at (66,-11.3) {2};
  \node[lab, anchor=west] at (69,-13.5) {\iflangja{後でまとめて反映}{applied later}};
\end{tikzpicture}
   (width ~116 mm)
\begin{tikzpicture}[x=1mm, y=1mm, font=\scriptsize\sffamily,
  >={Stealth[length=1.5mm]},
  lab/.style={font=\scriptsize\sffamily, align=center},
  num/.style={circle, draw=black, fill=white, text=black, line width=0.4pt,
              inner sep=0.4pt, font=\scriptsize\sffamily\bfseries}]
  % updates
  \node[draw, fill=white, align=center, inner sep=1.5pt] (up) at (39,22)
    {\iflangja{更新：ブロック 870，12，301}{updates to blocks 870, 12, 301}};
  \draw[->, line width=1.2pt, ln-accent] (up.south) -- node[num, right=0.8mm] {1} (39,8.2);
  \node[lab, anchor=west] at (43,14) {\iflangja{順に追記（シーケンシャル書き込み）}{appended in order (sequential writes)}};
  % disk strip
  \node[lab, anchor=east] at (17,4) {\iflangja{ディスク}{disk}};
  \draw[fill=white] (18,0) rectangle (116,8);
  \foreach \x/\b in {18/870, 32/12, 46/301} {
    \draw[fill=ln-accent!10] (\x,0) rectangle (\x+14,8);
    \node[font=\scriptsize\ttfamily] at (\x+7,4) {\b};
  }
  \foreach \x/\b in {66/12, 84/301, 102/870} {
    \draw[fill=ln-box] (\x,0) rectangle (\x+10,8);
    \node[font=\scriptsize\ttfamily] at (\x+5,4) {\b};
  }
  \node[lab, anchor=south west, text=ln-muted] at (18,8.3) {\iflangja{ジャーナル}{journal}};
  \node[lab, anchor=south, text=ln-muted] at (89,8.3) {\iflangja{本来の位置}{proper locations}};
  % apply later
  \draw[->, densely dashed] (25,0) .. controls (30,-15) and (100,-15) .. (107,0);
  \draw[->, densely dashed] (39,0) .. controls (42,-7) and (68,-7) .. (71,0);
  \draw[->, densely dashed] (53,0) .. controls (58,-10) and (84,-10) .. (89,0);
  \node[num] at (66,-11.3) {2};
  \node[lab, anchor=west] at (69,-13.5) {\iflangja{後でまとめて反映}{applied later}};
\end{tikzpicture}
   (width ~116 mm)
\begin{tikzpicture}[x=1mm, y=1mm, font=\scriptsize\sffamily,
  >={Stealth[length=1.5mm]},
  lab/.style={font=\scriptsize\sffamily, align=center},
  num/.style={circle, draw=black, fill=white, text=black, line width=0.4pt,
              inner sep=0.4pt, font=\scriptsize\sffamily\bfseries}]
  % updates
  \node[draw, fill=white, align=center, inner sep=1.5pt] (up) at (39,22)
    {\iflangja{更新：ブロック 870，12，301}{updates to blocks 870, 12, 301}};
  \draw[->, line width=1.2pt, ln-accent] (up.south) -- node[num, right=0.8mm] {1} (39,8.2);
  \node[lab, anchor=west] at (43,14) {\iflangja{順に追記（シーケンシャル書き込み）}{appended in order (sequential writes)}};
  % disk strip
  \node[lab, anchor=east] at (17,4) {\iflangja{ディスク}{disk}};
  \draw[fill=white] (18,0) rectangle (116,8);
  \foreach \x/\b in {18/870, 32/12, 46/301} {
    \draw[fill=ln-accent!10] (\x,0) rectangle (\x+14,8);
    \node[font=\scriptsize\ttfamily] at (\x+7,4) {\b};
  }
  \foreach \x/\b in {66/12, 84/301, 102/870} {
    \draw[fill=ln-box] (\x,0) rectangle (\x+10,8);
    \node[font=\scriptsize\ttfamily] at (\x+5,4) {\b};
  }
  \node[lab, anchor=south west, text=ln-muted] at (18,8.3) {\iflangja{ジャーナル}{journal}};
  \node[lab, anchor=south, text=ln-muted] at (89,8.3) {\iflangja{本来の位置}{proper locations}};
  % apply later
  \draw[->, densely dashed] (25,0) .. controls (30,-15) and (100,-15) .. (107,0);
  \draw[->, densely dashed] (39,0) .. controls (42,-7) and (68,-7) .. (71,0);
  \draw[->, densely dashed] (53,0) .. controls (58,-10) and (84,-10) .. (89,0);
  \node[num] at (66,-11.3) {2};
  \node[lab, anchor=west] at (69,-13.5) {\iflangja{後でまとめて反映}{applied later}};
\end{tikzpicture}

  \caption{ジャーナリング．散らばったブロックへの更新は，まず順にジャーナル
    へ追記され，後から本来の位置へ反映される．}
  \label{fig:l11-journal}
\end{figure}

\needspace{16\baselineskip}
\begin{keypoint}[章のまとめ]
  I/O デバイスはコマンド・データ・ステータスの各レジスタを見せ，PCIe などの
  インターコネクトを通じて到達され，オペレーティングシステムが標準化した
  インタフェースの背後でデバイスドライバによって操作される．デバイスは
  ブロックデバイス，キャラクタデバイス，ネットワークデバイスのいずれかであり，
  デバイスファイルとして表現される．CPU はメモリマップドI/O か I/Oポートで
  レジスタをアドレス指定し，割り込みかポーリングでデバイスのイベントを知り，
  プログラム制御I/O か DMA でデータを転送する．デバイスアクセスは通常
  カーネルを通るが，OS バイパスはデータ経路からカーネルを取り除ける．要求した
  スレッドはブロックするか，非同期に実行を続ける．ファイルへのアクセスは，
  ファイルシステム，汎用ブロック層，ドライバを通じてブロックデバイスに
  達する．仮想ファイルシステムは，ファイルディスクリプタ，dentry，iノード，
  スーパーブロックによってファイルシステムを統一する．ext2 はパーティションを
  ブロックグループに編成し，その iノードは直接ポインタと間接ポインタを使う．
  キャッシング，I/O スケジューリング，プリフェッチ，ジャーナリングが
  ディスクアクセスのコストを減らす．
\end{keypoint}

\end{document}
